\documentclass[11pt,a4paper,oneside,openany]{book}
\usepackage[xetex,top=24mm,bottom=25mm,left=27mm,right=25mm,headheight=15pt]{geometry}
\usepackage{fontspec,xeCJK}
\setmainfont{Times New Roman}
\setCJKmainfont[BoldFont=SimHei,ItalicFont=KaiTi]{SimSun}
\setCJKsansfont{Microsoft YaHei}
\setCJKmonofont{SimSun}
\usepackage{amsmath,amssymb,amsthm,mathtools,booktabs,longtable,array,graphicx,xcolor,fancyhdr}
\usepackage[unicode,bookmarks=false,colorlinks=true,linkcolor=blue!45!black,citecolor=blue!45!black,urlcolor=blue!45!black]{hyperref}
\newcommand{\Z}{\mathbb Z}\newcommand{\Q}{\mathbb Q}\newcommand{\F}{\mathbb F}\newcommand{\CC}{\mathbb C}
\newcommand{\SL}{\operatorname{SL}}\newcommand{\GL}{\operatorname{GL}}
\newcommand{\rad}{\operatorname{rad}}\newcommand{\diag}{\operatorname{diag}}
\newcommand{\src}[1]{\textsuperscript{\hyperref[source:#1]{[#1]}}}
\newtheorem{theorem}{定理}[chapter]
\newtheorem{lemma}[theorem]{引理}\newtheorem{proposition}[theorem]{命题}
\newtheorem{corollary}[theorem]{推论}
\theoremstyle{definition}\newtheorem{example}[theorem]{例}\newtheorem{conjecture}[theorem]{猜想}
\newcommand{\proved}{\textbf{已证}}
\newcommand{\conditional}{\textbf{在GRH下已证}}
\newcommand{\unfinished}{\textbf{尚未完成}}
\renewcommand{\proofname}{证明}\renewcommand{\contentsname}{目录}\renewcommand{\bibname}{参考文献}
\renewcommand{\chaptername}{讲}\renewcommand{\appendixname}{附录}\renewcommand{\tablename}{表}
\numberwithin{equation}{chapter}\setcounter{tocdepth}{1}
\setlength{\parindent}{2em}\setlength{\parskip}{4pt plus 1pt}\linespread{1.16}
\emergencystretch=3em\allowdisplaybreaks[2]\raggedbottom
\pagestyle{fancy}\fancyhf{}\fancyhead[L]{\small 任意元数平均：合数带来的新结构}\fancyhead[R]{\small\thepage}
\renewcommand{\headrulewidth}{0.3pt}
\fancypagestyle{plain}{\fancyhf{}\fancyfoot[C]{\thepage}\renewcommand{\headrulewidth}{0pt}}
\newif\iflectureappendix
\let\originalappendix\appendix
\renewcommand{\appendix}{\originalappendix\lectureappendixtrue}
\makeatletter
\def\@makechapterhead#1{\vspace*{8mm}{\parindent\z@\raggedright\normalfont
\ifnum\c@secnumdepth>\m@ne\if@mainmatter{\Large\sffamily\color{blue!40!black}\iflectureappendix\@chapapp\ \thechapter\else 第\,\thechapter\,讲\fi\par}\vskip12pt\fi\fi
\interlinepenalty\@M\LARGE\bfseries #1\par\nobreak\vskip24pt}}
\makeatother
\hypersetup{pdftitle={任意元数平均：合数带来的新结构与尚未完成的证明},pdfsubject={素数平均讲义续篇}}
\begin{document}\frontmatter
\begin{titlepage}\centering\vspace*{26mm}
{\Large 素数平均讲义续篇\par}\vspace{16mm}
{\Huge\bfseries 任意元数平均\par}\vspace{8mm}
{\LARGE 合数带来的新结构\par}\vspace{5mm}
{\Large 与尚未完成的证明\par}\vfill
\begin{minipage}{.8\textwidth}
从允许分母的多个素数出发，解释临界尺度、固定网络、因子转译、平方核心和GRH；把已证结果、条件性结果、局部资源和猜想分别接回原位置平均问题。
\end{minipage}\vfill
{\large 数学本科生研读版\quad 2026年9月\par}
\end{titlepage}


\chapter*{这是一份未完成研究的讲义}
\addcontentsline{toc}{chapter}{这是一份未完成研究的讲义}
前篇已经讲过素数元数的完整定理。现在把每次平均的数量改成任意整数 $k\ge2$。如果只是把字母 $p$ 换成 $k$，大部分矩阵仍然能写出来，但有些模逆元没了，有些块根本装不下，还有些位置数的G条件自动成立，却仍不一定能平均。

这一篇的目标是讲清这些新问题。我们不重证前篇的共同余数障碍、整数能量下降、CRT、本原列精确输送、Morris与Serre的深同余工具。用到时会明确说需要哪个结论、在新参数下怎样检查条件。新的证明内容会展开：非单位差分的配比、素因子指数给出的临界下界、固定网络和商均值模拟、平方及奇次幂核心，还有小整数生成单位群的数论。

\section*{四种状态，必须分别读}
\begin{longtable}{p{.20\textwidth}p{.71\textwidth}}
\toprule 标记&本篇中的含义\\\midrule
已证&已有文稿给出一般证明，或一般归约加完整有限证书。此处不等于再次独立审稿。\\
在GRH下已证&证明实际调用Dirichlet $L$ 函数的广义黎曼猜想；条件不能省略。\\
已有部分资源&例如已经构造一对可逆返回，但还没有证明完整根或终端覆盖。\\
猜想／待证&尚无完整全称证明。有限样本、固定模板失败均不能替代它。\\\bottomrule
\end{longtable}
本篇采用项目已有定理作为底稿，另外把其直接代入的小参数推论算清楚。旧文保留历史阶段的“尚未解决”，本篇不会照搬它们：例如十平均十六元后来已经有固定六步网络，平方族的全部入口也已经补齐。

\section*{从前篇带来的工具箱}
\begin{longtable}{p{.25\textwidth}p{.66\textwidth}}
\toprule 前篇结论&本篇怎样使用\\\midrule
本原中心化与 $G\mid n$&直接沿用，允许的素因子改为全部 $p\mid k$。\\
危险见证、重值、能量&只重讲非单位差分和偶数中点这两项合数改动。\\
真实返回与射影逆&继续逐项给位置计数和共同尺度，不重教矩阵记账。\\
根、局部化、深同余&沿用已证引理，另检查 $D>1$、行列式及危险模数。\\
精确本原列纤维&负责把有限同余结果提升到确切终点。\\
素数平均完整定理&用作合法商系统的求解器；模拟和最后块内清理在本篇证明。\\\bottomrule
\end{longtable}

\section*{建议的阅读次序}
第1—2讲先弄清研究目标和真正下界。第3—5讲是不依赖GRH的通用构造。第6讲把小元数结果放在正确位置。第7—9讲是目前最完整的新主线：平方元数如何把平均问题接到小单位群，再接到解析数论。第10—11讲解释奇次幂为什么分为已完成的2、3族与尚未完成的较大素数族。第12—14讲讲区间接续、失败路线和剩余任务。

每个来源标记如\src{S01}都指向末尾的来源表，压缩包内保留相应原文及相关核验资料。正文负责把数学讲通；原文层保留历史推导、具体证书和未采用的分支，便于后续研究接手。

\tableofcontents\mainmatter


\chapter{同一个G条件，现在有三个不同的临界问题}\label{ch:targets}
\section{允许的分母从一个素数变成一组素数}
本篇统一用 $k$ 表示每次平均的元数，$n$ 表示总位置数。非恒定有理输入照前篇中心化、本原化，得到整数零和 $X$，记差分最大公因数为 $G$。

共同余数障碍的证明不变，只需把“$\ell\ne p$”改为“$\ell\nmid k$”。因此一般必要条件是
\begin{equation}\label{eq:G}
G\mid n,\qquad \rad(G)\mid k.
\end{equation}
$\rad(G)$ 表示 $G$ 的不同素因子的乘积。第二项就是“$G$ 的每个素因子都整除 $k$”，并不要求 $G\mid k$。

例如六平均允许 $G=2^a3^b$，十平均允许 $G=2^a5^b$。同一个条件在不同位置数上会化简：六平均十元有 $G\mid10$，所以允许 $G=1,2$；六平均九元有 $G\mid9$，所以可能的 $G=1,3,9$ 全部被允许。这两个例子不是用了不同猜想。

\section{自动满足条件，不等于已经证明可解}
把命题“在 $n$ 个位置上，所有有理输入都恰由\eqref{eq:G}判定”记为 $\mathcal P(k,n)$。若 $\rad(n)\mid k$，由于 $G\mid n$，条件对每个输入自动成立。此时 $\mathcal P(k,n)$ 断言的就是\emph{所有输入都能平均}，仍然是一项需要证明的结论。

最直接的警示是24平均27元。27的素因子只有三，三整除24，所以G条件对所有输入都自动合法。然而
\[((-1)^{24},0,0,24)
\]
不能归零。下一讲会证明这个反例。可见“分母允许”远远不等于“位置上的操作足够”。

\section{M、H、N分别问哪件事}
排除一步全体平均的平凡维数 $n=k$，定义\src{S01}
\begin{align*}
M(k)&=\min\{n>k:\mathcal P(k,n)\},\\
H(k)&=\min\{n>k:\mathcal P(k,n),\ \rad(n)\nmid k\},\\
N(k)&=\min\{N>k:\text{每个 }n\ge N\text{ 都满足 }\mathcal P(k,n)\}.
\end{align*}
$M$ 问第一次在哪个位置数成功；$H$ 只看G条件真正能筛掉输入的维数；$N$ 则要求以后不再失败。

二平均有 $M=N=4$，但四元G自动合法，所以 $H=5$。六平均已有 $M=9,H=10$；仍未证明 $N=9$ 或 $N=10$。十平均同样有 $M=16,H=17$，最后阈值还没闭合。

不要从两个连续成功点推出后面全成功。若在零和配置中只平均一个子系统，它的均值一般不是零；冻结外部坐标再调用小维数解法，可能会到错误的目标。前篇的整体共同尺度约定在这里依然重要。

\section{合数真正改变的三件事}
第一，非零差分模 $k$ 未必可逆。例如模六的二不是单位。原来用 $i(a-b)$ 任意调余数的步骤失效，但这个公因子也可能提供新的整数混合比例。

第二，最早候选位置数可能远小于 $2k$。九平均的候选临界是十三，放不下两个九重块。平方族要改用 $(k,\sqrt k,1)$，不能沿用素数端点的 $(k,k,1)$。

第三，多个允许素数的分母变化必须由\emph{同一条实际路径}协调。CRT能拼接同余条件，却不能替一张超库存的表增加位置，也不能把两条各自有效的局部路径自动合并。

本篇后面的三项未完成任务正对应这些困难：较大素数奇次幂的临界核心、混合素因子的临界点、以及已解点之后的全部维数接续。\src{S02}\src{S03}


\chapter{临界尺度为什么要看素因子的指数}\label{ch:lower}
\section{只看最大素因子不够}
四平均六元有反例，八平均十二元也有反例。它们的底层素数都是二，却需要不同的剩余位置数量。决定这项下界的不是单独的 $p\mid k$，而是 $p$ 在 $k$ 中的指数。

写 $p^a\parallel k$，定义
\begin{equation}\label{eq:S}
s_p(k)=p^{\lceil a/2\rceil},\qquad S(k)=\max_{p^a\parallel k}s_p(k).
\end{equation}
$s_p(k)$ 是平方能被 $p^a$ 整除的最小正 $p$ 幂。例如
\[
S(4)=2,\quad S(8)=4,\quad S(9)=3,\quad
S(18)=3,\quad S(24)=4,\quad S(36)=3.
\]
为什么出现“半个指数”？下面的陷阱需要平均前后的赋值间隙都足够大，两边的要求交在这个平方根尺度上。\src{S01}

\section{一个不能逃走的最低赋值块}
记 $n=k+r$。考虑 $k$ 份相同值 $u$，其 $p$ 进赋值为 $\nu$；外部 $r$ 个值赋值至少为 $\nu+h$。取
\[m=\lceil a/2\rceil-1,\qquad h=a-m,\qquad h>m.
\]
一次非恒等操作若选了 $j$ 个外部位置，就用了 $k-j$ 份 $u$。只要 $1\le j<s_p(k)$，就有 $v_p(j)\le m<a$，并且
\[v_p(k-j)=v_p(j).
\]
重块部分的和有赋值 $\nu+v_p(j)$，外部部分的赋值更高，所以无法抵消。除以 $k$ 后，新均值的赋值恰为
\begin{equation}\label{eq:valuation-drop}
\nu'=\nu+v_p(j)-a<\nu.
\end{equation}
它又形成 $k$ 个同值位置，而留下的旧值与它至少相差 $a-v_p(j)\ge h$ 个 $p$ 因子。因此同样的状态结构重新出现。

\begin{theorem}[严格下界]
任意 $k\ge2$，在 $k<n<k+S(k)$ 都存在 $G=1$ 的不可达输入。在 $n=k+S(k)$，若 $\rad(n)\nmid k$，仍有这样的反例。
\end{theorem}
\begin{proof}
当 $2\le r<s_p(k)$，取
\[X_{k,r}=((-1)^k,0^{r-1},k).
\]
它零和、本原且含相差一的坐标，故 $G=1$。外部赋值至少 $a$，满足上述间隙。每个非恒等操作都按\eqref{eq:valuation-drop}重新产生非零的最低块，无法全零。
当 $r=s_p(k)=s$，唯一新选择是一次拿走全部外部位置。零和使所选和为 $-su$，所以输出为 $((-su/k)^k,u^s)$。本原代表是
\[((-1)^k,(k/s)^s),\qquad G=k/s+1.
\]
若它含不整除 $k$ 的素因子，则退出赋值陷阱后立即进入共同余数陷阱。因此仍不可达。$s\mid k$ 给 $\rad(k/s+1)\mid k\iff\rad(k+s)\mid k$。
最后处理 $r=1$。$k+1$ 个位置可达当且仅当某个输入已经等于总均值：第一次操作后为 $(a^k,-ka)$；若非零，下一步只有恒等或将参数换为 $-a/k$。零和向量 $(-k-1,-1,1^{k-2},4)$ 无零且 $G=1$，故给反例。
\end{proof}

\section{临界线B与非自动临界线C}
上述证明给出严格下界
\begin{equation}\label{eq:B}
B(k)=k+S(k)+\begin{cases}0,&\rad(k+S(k))\mid k,\\1,&\text{否则},\end{cases}
\qquad M(k),N(k)\ge B(k).
\end{equation}
而 $H$ 要跳过所有自动点，因此它的候选下界为
\begin{equation}\label{eq:C}
C(k)=\min\{n>k+S(k):\rad(n)\nmid k\},\qquad H(k)\ge C(k).
\end{equation}
不能总把 $C$ 写成 $B+1$：有时 $B$ 自己已经非自动，有时 $B$ 虽然由加一得到却仍自动。十平均与十二平均都是 $B=16$ 自动、$C=17$。

\begin{conjecture}[当前统一猜想]
任意 $k\ge2$ 有 $M(k)=B(k)$，并且 $N(k)=M(k)$。相应地预测 $H(k)=C(k)$。
\end{conjecture}
第一句问下界能否达到，第二句问成功能否延续。它们是两项尚未完成的量词；已经证明很多第一句的参数族，并不能替第二句。

\section{为什么少数边界没有加一}
达到最大尺度的素数唯一，记为 $p_*$。若其指数 $a_*\ge2$，则 $k/S$ 仍含 $p_*$，也含 $k$ 的其他全部素因子。因此 $k/S+1>1$ 与 $k$ 互素，边界必失败。
只有 $a_*=1$ 时，边界退出的 $G$ 才可能合法，此时必须
\[k/p_*+1=p_*^b,\qquad k=p_*(p_*^b-1).
\]
于是 $k+S=p_*^{b+1}$，成为允许素数幂的自动点。六平均九元、二十平均二十五元正属此形。

对素数幂 $k=p^a,a\ge2$，下界统一为
\[B(p^a)=p^a+p^{\lceil a/2\rceil}+1.
\]
前篇奇素数的 $2p+1$ 正是指数一的情况。24平均27元则有 $r=3<S(24)=4$，所以第一节的自动G反例确实落在赋值陷阱里。


\chapter{自动G点：什么时候真有一张固定网络}\label{ch:networks}
\section{所有有理输入可解，竟然能推出同一个词}
若固定一个平均词，对所有输入都达到全局均值，它的矩阵必为 $J_n/n$。条目属于 $\Z[1/k]$，因此 $n$ 的素因子都必须整除 $k$。这是固定网络的必要条件，24平均27元说明它还不充分。

更有用的是反向的一条有限维论证。若每个有理零和向量都可被某个词消去，取零和空间的一组有理基。先消去第一基向量；再为第二基向量的\emph{当前像}选一个消去词，接在后面；已经为零的像仍为零。有限次后，整个零和空间都被一个词消掉。平均词固定常值方向，所以这个词恰为 $J_n/n$。\src{S01}\src{S05}

因此
\[
\text{固定全输入网络存在}
\iff\bigl[\rad(n)\mid k\text{ 且 }\mathcal P(k,n)\bigr].
\]
右边并没有省掉 $\mathcal P(k,n)$；自动G合法仍须证明可达。

\section{六平均九元：四步里每一步在干什么}
先将任意六项平均为 $a$，另三项均值记为 $c$。第二步取三份 $a$ 与那三项，得到六份 $b=(a+c)/2$，另外三份仍为 $a$。
全局均值是 $\mu=(2a+c)/3$。第三步取两份 $a$、四份 $b$，均值恰为 $\mu$；剩下三项为 $a,b,b$，和为 $3\mu$。第四步取它们与三份 $\mu$，九项全等。\src{S09}

例如第一组均值 $a=3$、外部均值 $c=9$，第二步得 $b=6$，第三步得 $(2\cdot3+4\cdot6)/6=5$，恰为九元均值。剩余 $3,6,6$ 加三份五，也平均成五。整个词不依赖这些具体数值。

\begin{theorem}[四步网络族]\label{thm:four-network}
设 $k=ur$，$u\ge1$，且 $u+1\mid r$。在 $n=k+r$ 个位置上，存在固定四步 $k$ 平均网络。
\end{theorem}
\begin{proof}
第一步平均 $k$ 项为 $a$，余下 $r$ 项均值为 $c$。第二步取 $k-r$ 份 $a$ 与全部外部位置，得 $k$ 份 $b=((u-1)a+c)/u$，留下 $r$ 份 $a$。全局均值为 $\mu=(a+ub)/(u+1)$。
令 $h=ur/(u+1)=kr/(k+r)$，整除条件使它为整数。第三步取 $h$ 份 $a$、$k-h$ 份 $b$，均值为 $\mu$。剩余 $r$ 项总和必为 $r\mu$，第四步将它们与 $k-r$ 份 $\mu$ 平均即可。
\end{proof}
取 $k=m(m-1),r=m$，得到 $n=m^2$ 的网络：$2\to4,6\to9,12\to16,20\to25,30\to36,42\to49$ 等。若 $m$ 或 $m-1$ 为素数，它还与上一讲的下界相接，给出 $M(m(m-1))=m^2$。两者都合数时只能给一个上界，例如 $77\le M(72)\le81$。\src{S01}

固定全输入网络少于四步不可能：两步后，若新行纳入第一组外的位置，该输入位置的系数为 $1/k\ne1/n$；若没纳入，则支撑不足。未动行也不是全均值行。因此两步后没有任何一行正确，第三步又必留下 $n-k>0$ 行不动，不能全体正确。这是固定词下界，不是每个具体输入的最少步数。

\section{十平均十六元：用收缩误差配上最后的单点}
十六元不在上述四步族里，但有一张六步网络。更一般地，取
\[n=4^b,\qquad h=(n-1)/3,\qquad k=2h,\qquad b\ge1.
\]
先平均 $2h$ 个位置为 $a$，再平均 $h-1$ 份 $a$ 与全部其余 $h+1$ 个位置，得 $2h$ 份 $b_0$ 和 $h+1$ 份 $a$。留下一个 $a$ 作单点，其余整理为三个等大 $h$ 块，块值是 $(a,b_0,b_0)$。\src{S12}

现在每次合并两个 $h$ 块做一次 $k$ 平均，等于在三块值上做二平均。依次选块对 $(0,1),(1,2),(2,0)$ 循环。令 $a$ 的系数为一、$b_0$ 的系数为零，第 $j$ 次产生的新系数是
\[c_j=\frac{1-(-1/2)^j}{3},\qquad c_0=0,c_1=1/2.
\]
递推 $c_j=(c_{j-1}+c_{j-2})/2$ 可直接验证该公式。写 $m=2b$，在 $j=m-1$ 时
\[c_{m-1}=\frac{1+2^{1-m}}3=\frac{h+1}{3h+1},
\]
正好是全局均值里 $a$ 的系数。于是已有 $2h$ 份全局均值，最后将未达标的 $h+1$ 项与 $h-1$ 个正确位置平均，全部完成。总步数 $2+(m-1)+1=2b+2$。

对十平均十六元，后三次块混合依次得到 $(a+b_0)/2,(a+3b_0)/4,(3a+5b_0)/8$；最后一项就是六份 $a$ 与十份 $b_0$ 的总体均值。故 $M(10)=16$；结合十七元控制器，$H(10)=17$。六步没有被证明最短。

\section{张量网络：按坐标方向轮流平均}
若 $k=\prod p^{a_p}$，位置数 $B=\prod p^{b_p}$ 满足每个 $b_p\ge a_p$，把位置标为各素数进位坐标的笛卡尔积。每层选每个 $p$ 的 $a_p$ 条坐标轴，在对应的 $k$ 元纤维上平均。轮换轴，覆盖全部坐标，得到全体平均。\src{S05}

所需层数可取 $L=\max_p\lceil b_p/a_p\rceil$，操作数为 $LB/k$。证明只用不同坐标方向的全平均算子彼此交换。注意这是充分构造，不是所有固定网络的分类；六平均九元中 $6\nmid9$，仍有刚才的四步网络。

在 $3^b$ 个位置上，把九元六平均网络放到两条三进坐标轴组成的每个纤维，轮换覆盖所有轴，得到 $b\ge2$ 时的固定网络，长度
\[4\cdot3^{b-2}\lceil b/2\rceil.
\]
这些网络会作为下一讲块均值模拟的实际子程序。


\chapter{怎样真正把合数平均交给素数平均}\label{ch:transfer}
\section{为什么不能随时把六平均拆成二平均}
一个 $k$ 平均矩阵在 $n$ 维空间上的秩为 $n-k+1$。非空 $k$ 平均词的秩只会更小，所以不可能在\emph{整个空间}等于一个 $p<k$ 平均矩阵，其秩为 $n-p+1$。因此“六等于二乘三，所以能随时执行二平均”不成立。\src{S05}

但若每个商位置已有三份相同的值，两个商位置的六个实际副本恰好可以做一次六平均。这是在特殊子空间上模拟，完全可能。更一般地，$p\mid k$、每个商位置有 $c$ 份相同副本且 $k/p\mid c$ 时，一次商层 $p$ 平均用 $pc/k$ 次实际 $k$ 平均实现。

旧半块方法先制造这些相同副本，再调用二平均。奇数个商位置时，最后一半不能单独平均，便采用重叠入口：先平均 $2c$ 元组为 $a$，再取其中 $c$ 份与外部 $c$ 项平均为 $b=(a+v)/2$，留下 $c$ 份 $a$，形成商值 $(a,b,b,\ldots)$。在危险奇素数处，$a=b$ 等价于 $a=v$，所以见证被保留。\src{S05}

\section{后来的简化：起初块内不相等，也照样能模拟}
其实不必先让每块相等。把 $n$ 个位置分成等大的 $c$ 元块，只记录每块的\emph{均值} $z_i$。若选 $p$ 个块，对其全部 $pc$ 个位置执行一个固定全平均网络，它们的新均值就是
\[\frac1{pc}\sum_{\text{这些位置}}x_j=\frac1p\sum_{\text{这些块}}z_i.
\]
未选块均值不变。因此商均值严格跟随一条 $p$ 平均路径，被选块内部还顺便变得全等。\src{S06}

商均值最后都正确，还可能有某块一直没选过、内部不等。补一轮覆盖所有块的 $p$ 元块组即可：此时每个块均值相同，再平均不会改变目标，却会清掉所有块内差异。这一清理步骤不能省。

\section{先把危险素数从块均值里保护下来}
我们还须选择分块，使商输入合法。已有同时好划分引理说：块数 $m\ge5$，每个块大小在所有需保护素数处可逆，而且危险素数个数不超过 $m-1$，就能选到各素数处块均值不全相同的划分。\src{S29}

它的直观证明是逐个修复素数：交换计划中两个块的一个原位置，会改变两个块均值。每个已经保护的素数至多禁止一条块间边；删去少于 $m-1$ 条边的完全图仍连通。如果允许边两侧都找不到不同余位置，原输入就会全同余，矛盾。交换只是修改分组计划，最终仍对原位置实施平均。

\begin{theorem}[素数因子转译]\label{thm:factor}
取 $p\mid k$ 为素数，假设 $k/p\mid n$。定义
\[
c=p^{v_p(k)-1}\prod_{\ell\mid k,\ell\ne p}\ell^{v_\ell(n)},\qquad m=n/c.
\]
若 $m=p$，或 $p=2,m\ge4$，或 $p$ 奇且 $m\ge2p+1$，则 $\mathcal P(k,n)$ 成立。
\end{theorem}
\begin{proof}
$c\mid n$，$k\mid pc$，且 $pc$ 与 $k$ 有相同素因子支持。上一讲的张量网络可在任意 $pc$ 个位置实现全平均。
商维数 $m$ 的非 $p$ 素因子都不整除 $k$；它们是原输入已有的危险素数，并且 $c$ 在这些地方可逆。当 $m\ge5$ 时，用同时好划分安排 $m$ 个 $c$ 块；$m=4,p=2$ 时没有危险奇素数；$m=p$ 时直接对全部位置使用网络。
对合法商输入调用前篇素数定理，用 $pc$ 元网络逐步模拟其平均。所有块均值相同后，按本节开头的覆盖清理完成所有原位置。
\end{proof}

例如十五平均33元，取 $p=5,c=3,m=11$。每次选五个三元块，直接做一次十五平均；商层五平均十一元已解。十五平均35元则取 $p=3,c=5,m=7$，调用三平均七元。两者都是原来就有的位置，没有添加副本。

六平均14元取 $p=3,c=2,m=7$，六平均22元取 $p=3,c=2,m=11$，也都直接完成。这两项说明旧六平均未解表需要更新。

\section{哪些整除族已经全部解决}
\begin{corollary}
任意偶数 $k$ 的全部 $k$ 倍数维数均满足G判据。任意 $k\ge2$ 的全部 $n=ak$、$a\ge3$ 也满足G判据。
\end{corollary}
\begin{proof}
偶数取 $p=2$，商 $m$ 为偶数且至少二，故满足转译条件。
设 $k$ 奇，将倍数 $a=uv$ 分解为 $u$ 只含 $k$ 的素因子、$\gcd(v,k)=1$。对 $p\mid k$，商为 $m=p^{v_p(u)+1}v$。$v=1$ 时是 $p$ 或至少 $p^2$；$v\ge3$ 时至少 $3p$；$v=2$ 时因 $a\ge3$ 有 $u>1$，取 $p\mid u$，商至少 $2p^2$。均满足前篇阈值。
\end{proof}
更具体地，四平均覆盖全部偶数 $n\ge8$，另有 $n=4$；六平均覆盖全部三的倍数 $n\ge6$。后一结论中，去掉所有三因子后的商为二或至少四时用半块转译，若商为一，则用上一讲的 $3^b$ 固定网络。\src{S05}

这些都是整除族，不是连续区间。不能把“全部 $n=ak$、$a\ge3$”读成“全部 $n\ge3k$”。剩下不整除的邻点，正是接续问题。


\chapter{整数下降怎样跨过合数障碍}\label{ch:tail}
\section{非单位差分也能帮忙}
前篇常用差分 $a-b$ 的模逆来调配比。现在已有至少 $k$ 份 $a$ 和至少 $k$ 份 $b$，且 $a\ne b$。对合数 $k$，设 $g=\gcd(a-b,k)>1$，模逆不存在，但只取 $k/g$ 份 $a$、其余取 $b$ 就得到
\begin{equation}\label{eq:nonunit}
\frac{(k/g)a+(k-k/g)b}{k}=b+\frac{a-b}{g}\in\Z.
\end{equation}
两种值都留有未选副本，因此这一步安全。\src{S04}

例如六平均中 $a=5,b=1$，差分四与六的最大公因数为二。取三份五、三份一，得到六份三。不可逆的差分没有让我们无路可走，反而给出了一个只用两种值的整数平均。

\section{偶数元数唯一新增的碰撞：两个配比同为中点}
把两种配比对调，两个新均值是
\[c_1=b+(a-b)/g,\qquad c_2=a+(b-a)/g.
\]
其差为 $(2/g-1)(a-b)$，所以只在 $g=2$ 时相等，此时都是中点。

若恰有三个重值 $a<b<c$，原本想用两种配比之一避开第三值。唯一可能同时撞上第三值的情形，是两端平均到中间：
\[a+c=2b,\qquad\gcd(c-a,k)=2.
\]
此时相邻差分必须是单位；若相邻差分也非单位，其两值混合的均值严格在相邻两值之间，不会等于第三个。于是又可以用前篇的大小顺序论证：有轻值就保留另一极端，极端有富余就只取它一份；最后只剩相反值加零的终端。

\section{其余分类怎样接回前篇}
大背景值的证明分两支：差分非单位时直接用\eqref{eq:nonunit}；差分单位时仍可唯一选择 $i$，使 $(a^i,b^{k-1-i},x)$ 为整数平均。若所有选择碰撞，所有值被迫在同一有限等差列中，前篇的系数跨越打包继续成立。

需要补的一点是，等差列中的 $x-b=-j(a-b)$ 也可能是非单位。遇到它不能继续当作可逆差分；再次使用\eqref{eq:nonunit}即可。三个以上重值用上一节的中点分析，恰两个重值则各留 $k$ 份，再保护 $d_k(n)$ 个见证，在剩余池应用一般EGZ。

\begin{lemma}[一般EGZ由素数版本推出]
任意 $2k-1$ 个整数中，有 $k$ 项的和被 $k$ 整除。
\end{lemma}
\begin{proof}
素数情形已在前篇证明。若 $k=uv$，从 $2uv-1$ 项中反复取出 $2v-1$ 个不交的 $u$ 元组，使每组和被 $u$ 整除；最后一次取组前仍至少有 $2u-1$ 项。将各组和除以 $u$，应用 $v$ 的版本，选 $v$ 组的并集。对因子个数归纳。
\end{proof}

\begin{theorem}[任意元数的已证尾部]
对 $k\ge3$，若
\[n\ge4k-1+d_k(n),\qquad d_k(n)=\#\{\ell\mid n:\ell\nmid k\},
\]
则 $\mathcal P(k,n)$ 成立。特别
\[N(k)\le4k+\lceil\log_2k\rceil+1\le5k.
\]
\end{theorem}
\begin{proof}
从 $kX$ 初始化两个重值。上面已核对前篇分类的全部合数改动。恰两重值且都不足 $2k-1$ 份时，各留 $k$ 份，再保护至多 $d_k(n)$ 个位置，余池至少 $2k-1$ 项，且每个实际值不足 $k$ 份；一般EGZ给非恒等安全整数平均。
因此每个非终端都有保持储备的整数续步，平方和严格下降，有限步进入相反值终端。最后的统一界用 $d_k(n)\le\lfloor\log_2n\rfloor$ 及整数函数 $n-\lfloor\log_2n\rfloor$ 不减，与前篇相同。
\end{proof}

这项结论已经保证每个 $k$ 的最终阈值\emph{存在}。未完成的是精确的最优阈值。旧文曾说合数还不知道是否有最终阈值，那是早于本定理的阶段。

\section{精度与算法：这里究竟保证多快}
对合法整数零和 $X$，在 $kX$ 上做整数下降，原 $X$ 的中间值都属于 $(1/k)\Z$。合数情况下，分母可以是 $k$ 的真因子，例如六平均可出现二分之一、三分之一；不能沿用“分母只有一或 $p$”的素数断言。

整数能量每次至少降二，给粗界 $T\le(k^2/2)\sum X_i^2$。这是关于数值大小的界，不是输入位长的多项式界。已知充分范围内的\emph{判定}则很快：算 $G$，反复除去 $\gcd(G,k)$，直到得到一或剩余与 $k$ 互素；不需要分解巨大 $G$。

四平均已有 $N(4)\le16$，六平均已有 $N(6)\le24$。还有一个方便终端：对这两种元数，任意维数一旦产生 $k$ 个零就能完成。非零支持不超过 $k$ 时直接一次平均；其他情况四平均保留一个零后用偶数维数转译，六平均保留 $n\bmod3$ 个零后用三倍数维数定理。含零子系统自身零和且本原 $G=1$，调用合法。\src{S04}


\chapter{先看几个小合数：六平均为什么值得单独讲}\label{ch:small}
\section{两个新问题：用哪个模数，能缩放哪些倍数}
前篇已经解释过，把许多位置压成三种重数后，可以用二乘二矩阵记录返回。合数平均继续使用这一工具，但有两件事要重新确定。第一，总位置数的某些素因子已经允许出现在分母里，它们不该再成为障碍。第二，已经做出了加法式的剪切，还不一定做得出完成归约所需的乘法式缩放。

六平均十元清楚地展示第一件事，十平均与十二平均十七元展示第二件事。它们都已经完成，适合当作一般构造的具体模型。四平均七元、八平均十三元和九平均十三元的意义，则是提前展示后面会出现的参数族。

\section{六平均十元：允许偶数的差分}
十元本原零和输入满足 $G\mid10$。六平均允许分母二和三，所以合法的 $G$ 是一或二；五和十不合法。\src{S10}

例如 $(1^5,(-1)^5)$ 的 $G=2$。先选三份一、三份负一，产生六个零；再选剩下四个非零位置和两个零，十项便全为零。这个两步例子直观说明：排除偶数 $G$ 会排除真实可解的输入。与之相对，$(1^9,-9)$ 的 $G=10$，模五始终全为一，确实不可解。

先留四个位置，平均另外六个，得到六份 $a$。保留位置中可以安排两个模五不同的值，于是能选一个 $w\not\equiv a\pmod5$。再平均三份 $a$ 和其余三个保留值，就进入
\[
K(u,v)=(u^6,v^3,-6u-3v).
\]
总和为零给出 $u=(-3a-w)/6$、$v=a$，因此
\[u-v\equiv(a-w)/6\ne0\pmod5.
\]
取 $x=u,y=u-v$，参数清分母并本原化后，
\[G(K)=\gcd(y,10).
\]
我们要搬动的列向量满足 $\gcd(x,y)=1$ 和 $5\nmid y$。它完全可能有偶数 $y$。若强迫 $y$ 在模十下可逆，就会误删全部 $G=2$ 的合法输入。因此这里的同余群应该保留模五条件。

\section{一个两步返回补齐了正向逆}
三个一步返回分别选择 $(u^3,v^3)$、$(u^5,w)$、$(u^3,v^2,w)$，其中 $w=-6u-3v$。它们在差坐标上是
\[
A=\frac12\begin{pmatrix}2&-1\\0&-1\end{pmatrix},\quad
B=\frac16\begin{pmatrix}-4&3\\-10&9\end{pmatrix},\quad
C=\frac16\begin{pmatrix}-4&1\\-10&1\end{pmatrix}.
\]
再做一个两步返回：先平均 $u^4,v,w$，得到六份 $t=-(u+v)/3$；然后平均三份 $t$、一份 $u$、两份 $v$。留下三份 $t$ 和一份 $u$，所以
\[
R=\frac16\begin{pmatrix}1&-1\\5&-3\end{pmatrix}.
\]
矩阵计算给出
\[
R^4=-\frac1{324}I,\qquad(AR)^2=\frac1{36}I,
\qquad(RAC)^2=\frac1{216}I.
\]
这些等式的用途与前篇相同：把一个因子移到等式另一侧，就把它的射影逆写成只向前执行的返回。例如 $A^{-1}\sim RAR$。这样，$A,C,R$ 都有正向逆，进而
\[\Sigma=C^{-1}B=\begin{pmatrix}1&-1\\0&-1\end{pmatrix}
\]
也能执行。注意六重块和三重块大小不同，$\Sigma$ 的来源是这一返回词。

记 $U(z),L(z)$ 为上、下剪切，$\Delta(z)=\diag(1,z)$。因为 $\Sigma A=\Delta(1/2)$，我们有 $D=\Delta(2)$ 及其逆，从而
\[\Sigma D\Sigma D^{-1}=U(1/2).
\]
整数幂与共轭给出所有二进上剪切。接下来算
\[
F=U(-1)C=\begin{pmatrix}1&0\\-5/3&1/6\end{pmatrix},\qquad
DFD^{-1}F^{-1}=L(-5/3).
\]
将后者的逆执行三次，得到 $L(5)$。原来需要的模五条件，恰好出现在下剪切的系数里。进一步由 $L(5/3)F=\Delta(1/6)$ 得到三的缩放，最后可以执行
\[U(\Z[1/6]),\qquad L(5\Z[1/6]),\qquad
\Delta(\Z[1/6]^\times).
\]

\section{把任意合法输入送到一个看得见的终点}
前篇的陪集回路方法此时只需处理 $\mathbf P^1(\F_5)$ 的六个点。代表取 $\infty,0,\pm1,\pm2$；闭合回路的主元都是 $\pm1,\pm2$，已经可以缩放。于是整个 $\Gamma_0(5)$ 都有正向实现。

对合法本原列 $(x,y)$，求 $\alpha x+\beta y=1$，再调整系数使 $5\mid\alpha$。这是可行的，因为调整 $\alpha\mapsto\alpha+jy$ 时，$y$ 在模五下可逆。于是
\[
\begin{pmatrix}y&-x\\\alpha&\beta\end{pmatrix}
\binom{x}{y}=\binom01
\]
给出所需归约。核心变成 $(0^6,v^3,-3v)$，再选三个 $v$、一个 $-3v$ 和两个零做一次六平均，就结束了。

这证明 $H(6)=10$。九元的固定网络则证明 $M(6)=9$。两个结果回答不同的问题：九元研究如何同时混合全部输入，十元研究同余障碍是否是唯一障碍。

\section{十平均、十二平均十七元：剪切够了，缩放还不够}
两种核心是 $(u^{10},v^6,w)$ 和 $(u^{12},v^4,w)$。它们共用两步入口；因为十七是唯一危险素数，合法域都是 $17\nmid y$。\src{S11}

两步返回的物理形式也相同：平均指定位置得到 $k$ 份 $t$，保留 $r$ 份 $t$ 和一个未参加第一步的旧值 $z$，再平均其余 $k$ 个位置。由总和零，输出是
\[u'=-(rt+z)/k,\qquad v'=t,\qquad w'=z.
\]
十平均的一组周期返回为
\[
R=\frac1{50}\begin{pmatrix}16&-6\\51&-16\end{pmatrix},
\quad E=\frac1{50}\begin{pmatrix}-35&6\\-85&16\end{pmatrix},
\quad R^2=-I/50,\quad(RE)^2=I/2500.
\]
十二平均的一组为
\[
R=\frac1{36}\begin{pmatrix}2&-2\\17&-8\end{pmatrix},
\qquad R^4=-I/5184.
\]
配合其他几个同型返回，二者都得到 $\Delta(2)^{\pm1}$、$\Sigma$、$U(\Z[1/2])$ 和 $L(17\Z[1/2])$。原稿保存全部选取表及推导，附录的来源包一并保留。

然而模十七的带符号二幂只占八种非零余数，全部非零余数有十六种。要继续照搬陪集回路证明，还需要新倍率。十平均从返回比值得到 $\Delta(5)$，并先构造 $\Delta(25)^{-1}$，再用
\[\Delta(5)^{-1}=\Delta(25)^{-1}\Delta(5)
\]
得到它的逆。十二平均以同样思路补齐三的缩放。最终工作环分别是 $\Z[1/10]$ 和 $\Z[1/6]$。

现在可选模十七单位代表
\[
\{\pm1,\pm2,\pm4,\pm5,\pm8,\pm10,\pm20,\pm40\},
\]
或
\[
\{\pm1,\pm2,\pm3,\pm4,\pm6,\pm8,\pm12,\pm24\}.
\]
每套都覆盖全部十六种余数，并且代表本身是可执行的单位。前篇的回路编译与贝祖输送于是完成证明：$H(10)=H(12)=17$。固定网络另外给出 $M(10)=M(12)=16$。

\section{小例子怎样汇入一般族}
四平均七元对应下一讲的 $t=2$，九平均十三元对应 $t=3$。八平均十三元则对应后面的 $k=2s^2,s=2$。它们已经有无条件证明；一般族要解决的问题，是怎样在每个参数处都找到足够的可逆倍率。

这些例子也说明，发现程序与最终论证可以各有明确位置。程序先找到少量返回，证明随后只使用显示的矩阵恒等式、位置选取和完整群覆盖。原始运行中的千步样本说明编译器能输出实际路径，却不能说明路径最短；后面的参数族也尚未给出统一的最短词算法。


\chapter{平方元数：一套可以写成公式的控制器}\label{ch:square}
\section{为什么尝试平方}
令 $k=t^2$，总位置数取 $n=t^2+t+1$，其中 $t\ge2$。核心形状是
\[K(u,v)=(u^{t^2},v^t,w),\qquad w=-t^2u-tv.
\]
平方给出一个自然的库存比例：大块恰好能分成 $t$ 组，每组 $t$ 个位置。后面会看到，这个比例还让某个返回矩阵满足三次周期。

这一讲先解决已经进入核心以后的控制问题，下一讲再解决任意输入的入口和最后的算术条件。这样可以清楚看到：哪些是矩阵恒等式，哪些还需要关于整数的定理。\src{S13}

始终取差坐标 $x=u,y=u-v$。本原化参数后，$G(K)=\gcd(y,n)$，而 $\gcd(n,t)=1$，所以合法核心恰好满足 $\gcd(y,n)=1$。

\section{三个一步返回}
分别平均
\[(u^{t^2-t},v^t),\qquad(u^{t^2-1},w),\qquad
(u^{t^2-t},v^{t-1},w).
\]
第一种留下 $t$ 份旧 $u$，第二种留下 $t$ 份旧 $v$，第三种留下 $t$ 份旧 $u$；其余单点也明确保留。三个实际作用为
\[
\mathcal A=\frac1t A,\quad A=\begin{pmatrix}t&-1\\0&-1\end{pmatrix},\qquad
\mathcal B=\frac1{t^2}B,\quad B=\begin{pmatrix}-t-1&t\\-n&t^2+t\end{pmatrix},
\]
\[
\mathcal C=\frac1{t^2}C,\qquad C=\begin{pmatrix}-t-1&1\\-n&1\end{pmatrix}.
\]
例如第三种的新大块均值为 $-(tu+v)/t^2$，新小块值为 $u$。相减后便得到 $C$ 的第二行。这说明矩阵并非猜出来的：每个系数都能从选取的位置数直接读出。

\section{周期把不可逆平均变成可逆的方向控制}
矩阵 $C$ 的迹是 $-t$、行列式是 $t^2$，特征恒等式给出
\[C^2+tC+t^2I=0,\qquad C^3=t^3I.
\]
于是实际返回满足 $\mathcal C^3=t^{-3}I$。连续三次返回缩小了状态，却把方向完全还原。这正是可以实现射影逆的原因。

另外，$\operatorname{tr}(CA)=0$，从而
\[(\mathcal C\mathcal A)^2=t^{-3}I.
\]
因此 $\mathcal A$ 也有正向射影逆。最后
\[C^2B=t^3\Sigma,\qquad
\Sigma=\begin{pmatrix}1&-1\\0&-1\end{pmatrix}.
\]
至此，$\mathcal A,\mathcal B,\mathcal C,\Sigma$ 的射影逆都有真正的向前执行方式。

因为 $\Sigma\mathcal A=\Delta(1/t)$，我们已有 $\Delta(t)^{\pm1}$。交换子给出
\[\Sigma\Delta(t)\Sigma\Delta(t)^{-1}=U((t-1)/t).
\]
取幂和共轭后，得到 $U((t-1)R_0)$，其中 $R_0=\Z[1/t]$。借助 $C$ 的横向作用，前篇的同余子群工具给出一个深层
\[\Gamma((t-1)^2n^3,R_0).
\]
这里先出现的是带因子 $t-1$ 的剪切。这足以准备高精度，却还没有得到任意大小的上剪切。

\section{两步返回供应所有小整数倍率}
选 $0\le j\le t-1$，先平均
\[(u^{t^2-j-1},v^j,w),
\]
得到 $t^2$ 份新值。保留其中 $t$ 份，以及一个未选的旧 $u$，再平均其余 $t^2$ 个位置。实际矩阵是 $M_j/t^3$，其中
\[M_j=\begin{pmatrix}1&j-t\\n&(t+1)(j-t)\end{pmatrix}.
\]
所有选择都在库存内：第一次之后至少剩一份旧 $u$，新值则足够保留 $t$ 份。

最有用的地方是这些矩阵第一列相同。因而
\[M_0^{-1}M_j=\Delta((t-j)/t).
\]
$\det M_0=t^3$，是 $R_0$ 中的单位；它在相应有限商中保持已取得的深层。应用前篇的有限商幂论证，可以用正向词实现 $M_0$ 的射影逆。配合 $\Delta(t)$，便得到
\[\Delta(d),\qquad 1\le d\le t.
\]
这一步首先给正向倍率。某个 $d$ 若含有 $n$ 的素因子，其逆会破坏合法域，所以不能一概立即使用。

\section{为什么必须专门处理除以二}
我们希望把旧剪切中的 $t-1$ 除掉。单纯把更多素数加入环并不总能做到，因为 $t-1$ 可能还含不可逆因子。有效办法是重新计算一个交换子：一旦 $\Delta(2)$ 和逆都可执行，就有
\[\Sigma\Delta(2)\Sigma\Delta(1/2)=U(1/2).
\]
新的系数已经不带 $t-1$。这是新恒等式提供的改进。

若 $t$ 为偶数，二已经是 $R_0$ 的单位，初始深层的水平也是奇数，可以用有限商恢复倍率的逆。若 $t$ 为奇数，二可能碰到旧水平中的 $t-1$，要用一个专门的行列式一矩阵：
\[
H=\frac1{2t}\Delta(2)C\Delta(2)
=\begin{pmatrix}
-(t+1)/(2t)&1/t\\-n/t&2/t
\end{pmatrix}\in\SL_2(R_0).
\]
它属于 $\SL_2(R_0)$ 的关键是 $t+1$ 为偶数，且行列式等于一。$H$ 本身由已实现因子射影给出；它的某个正整数次幂落入深层，所以有正向射影逆。再用
\[C\Delta(2)H^{-1}=2t\Delta(1/2)
\]
就实现了所需的除二倍率。

于是，令 $R=\Z[1/(2t)]$，我们得到 $U(R)$。下剪切来自另一个直接恒等式：
\[U(-1)\mathcal C\Delta(t^2)=L(-n/t^2).
\]
取幂、共轭，得到全部 $L(nR)$。平方控制器的主要工作至此完成：原位置平均已经供应了一整套上剪切、$n$ 倍下剪切和若干倍率。

\section{继续扩大允许倍率的范围}
对于 $1\le d\le t$ 且 $\gcd(d,n)=1$ 的整数，$\Delta(d)$ 已有正向实现。现在用新的同余层，可以恢复它的逆，进而把 $d$ 的素因子加入局部化环。记所得环为 $R_t$。

这里扩大的是\emph{核心方向矩阵}的系数环。真实坐标仍只有 $k$ 所允许的分母。例如 $\Delta(1/5)$ 与 $\diag(5,1)$ 表示同一个射影作用；实际返回还带一个由平均词决定的共同尺度。因而方向矩阵出现五的分母，不表示许可操作突然变成了五平均。这正是前篇共同尺度约定在此处的用途。

可逆倍率越来越多，模 $n$ 能改变的余数也越来越多。不过“很多”还不等于“全部”。下一讲会把唯一剩下的乘法问题写成一个有限群等式。至于本讲的周期、除二和剪切公式，都对每个 $t\ge2$ 成立，不依赖 GRH。


\chapter{从任意输入到平方核心，再到一个有限乘法群}\label{ch:entry}
\section{先把所有同余见证集中到一个单点}
上一讲假设已经有 $(t^2,t,1)$ 核心。现在说明任何合法输入都能进入它。这部分已经对所有 $t\ge2$ 证明，而且不依赖后面的单位群结论。\src{S14}\src{S15}

记 $n=t^2+t+1$，$d=\omega(n)$ 为不同素因子个数。因为
\[2^{t+1}>t^2+t+1,
\]
必有 $d\le t$。固定一个位置，为每个 $\ell\mid n$ 再保护一个与它模 $\ell$ 不同的位置。总共至多保护 $t+1$ 个位置，余下足够平均一次，得到
\[(a^{t^2},c_1,\ldots,c_{t+1}).
\]
令 $\delta_i=c_i-a$。每个危险素数处，至少有一个 $\delta_i$ 非零。问题是：能不能把这些可能分散在不同单点上的非零信息集中起来？

前篇的块与单点交换只需一个大块。平均 $t^2-1$ 份 $a$ 和 $c_i$，将剩下的一份旧 $a$ 登记为新单点。差分作用为
\[A_i=I+(\lambda\mathbf1-e_i)e_i^T,\qquad\lambda=-1/t^2.
\]
三个不同角色 $i,j,h$ 给出交换子
\[[A_h^{-1}A_jA_h,A_i]
=I+\frac{t^2-1}{t^2}(e_h-e_j)e_i^T.
\]
只为调整模 $\rad(n)$ 的信息时，形式逆可由有限群中的正整数幂替代。得到的实际数值未必是有理逆的数值，但余数恰好相同，足以完成这一入口任务。

这个交换子能给 $\delta_j$ 加上 $\delta_i$ 的倍数，唯一可能失效的素数由
\[\gcd(t^2-1,n)=\gcd(t+2,3)
\]
控制。因此至多只有三失效。若三失效，先把一个模三非零的角色选作目标 $j$。其他危险素数逐个处理：在当前素数处选非零的 $\delta_i$，用CRT选择交换子重复次数，让这一次只改变当前素数处的目标余数。第三个单点提供缓冲。最后同一个 $\delta_j$ 在全部危险素数处非零。

留下目标单点 $w=c_j$ 和 $t$ 份旧 $a$，平均其余 $t^2$ 项。所得 $u,v$ 满足
\[v=a,\qquad u=(-ta-w)/t^2,\qquad
u-v\equiv(a-w)/t^2\pmod\ell.
\]
右边在每个 $\ell\mid n$ 处非零，于是到达合法平方核心。

这里的复用值得注意：旧证明在 $n>2k$ 中使用这一交换子，是因为后面还要建第二个大块。交换子本身只需要三个单点。把它的真实前提分离出来，就能用于 $n=t^2+t+1<2t^2$ 的新范围。

\section{剩下的是哪个有限群}
定义
\begin{equation}\label{eq:Ht}
H_t=\langle-1,\ d\bmod n:\ 1\le d\le t,\ (d,n)=1\rangle
\subseteq(\Z/n\Z)^\times.
\end{equation}
这个群只记录上一讲已经可以正向、可逆实现的倍率。局部化环 $R_t$ 的单位在模 $n$ 下的像恰好是 $H_t$。

完整上剪切和 $n$ 倍下剪切，结合前篇的同余工具，给出
\[\langle U(R_t),L(nR_t)\rangle=\Gamma_1(n,R_t),
\]
其中 $\Gamma_1$ 中矩阵的左下项为零模 $n$，两个对角项为一模 $n$。它能细调当前列，却不能任意改变第二坐标的模 $n$ 余数。改变这个余数正是倍率的工作。

\begin{proposition}[平方族的充分算术接口]\label{prop:Ht}
若 $H_t=(\Z/n\Z)^\times$，则 $t^2$ 平均在 $n=t^2+t+1$ 个位置上的完整G判据成立。
\end{proposition}
\begin{proof}
先用上一节进入合法核心，将参数取本原整数代表 $(x,y)$。由 $\gcd(y,n)=1$，可求
\[\alpha x+\beta y=1,\qquad n\mid\alpha.
\]
令 $F=\left(\begin{smallmatrix}y&-x\\\alpha&\beta\end{smallmatrix}\right)$，则 $F(x,y)^T=(0,1)^T$。

由 $H_t$ 满群，可以找到 $R_t$ 的单位 $\epsilon\equiv y\pmod n$。矩阵
\[D_\epsilon F\in\Gamma_1(n,R_t),\qquad
D_\epsilon=\diag(\epsilon^{-1},\epsilon)
\]
的每一项同余条件都成立：左上项为 $y/\epsilon\equiv1$，右下项由行列式一也为一，左下项仍被 $n$ 整除。$D_\epsilon$ 射影等价于已可执行的 $\Delta(\epsilon^2)$，故 $F$ 可以正向实现。

输出方向 $(0,1)$ 对应 $(0^{t^2},(-1)^t,t)$。选全部 $t$ 份负一、单点 $t$ 及 $t^2-t-1$ 个零，做一次平均即全零。
\end{proof}

这个命题把“还差一点算术”变成了可独立研究的有限群问题。其正向实现、任意输入入口和最终操作都已明确。

\section{为什么换更多同形终点并没有绕过它}
如果 $H_t$ 不满，能否换一个终端方向？先看所有在本核心内一次选出零和组的方向。

选 $b$ 份 $v$，不选单点时，零和方程为 $t^2x-by=0$；本原方向的第二坐标是
\[y=\frac{t^2}{\gcd(b,t^2)},\qquad0\le b\le t.
\]
若选单点，方程为 $-(t+1)x+(t-b)y=0$，第二坐标是
\[y=\frac{t+1}{\gcd(t-b,t+1)}.
\]
两种 $y$ 的模 $n$ 余数都属于 $H_t$。第一种只用不超过 $t$ 的因子；第二种还用到
\[t+1\equiv-t^2\pmod n.
\]
边界 $b=0$ 或 $b=t$ 直接给 $y=1$。

所以继续枚举这种一步零和方向，不能补齐 $H_t$ 外的尺度类。这个限制针对当前核心和当前倍率菜单。它没有排除离开核心再返回、用更长终端词，或通过状态相关约分得到新作用。

\section{一个等式把下一步交给数论}
命题\ref{prop:Ht}只要求证明有限交换群被一批小整数生成。它不再涉及平均的选点过程，已经能直接调用数论工具。下一讲会看到：素数模数有初等证明；一般模数有无条件的大参数结果；在GRH下可以把剩余有限范围接齐。


\chapter{GRH怎样进入一个平均问题}\label{ch:characters}
\section{用特征检测一个子群是否太小}
我们要证明 $H_t$ 是模 $n=t^2+t+1$ 的全部单位群。有限交换群有一个方便的检测器：如果一个子群 $H$ 真小于群 $\mathcal U$，商群 $\mathcal U/H$ 非平凡，因而存在一个非平凡乘法同态
\[\chi:\mathcal U\longrightarrow\{z\in\CC:|z|=1\}
\]
在 $H$ 上恒等于一。将它在非单位处定义为零，就得到一个Dirichlet特征。

因此，$H_t$ 不满等价于存在非主特征满足
\[\chi(-1)=1,\qquad\chi(d)=1\quad(1\le d\le t,\ (d,n)=1).
\]
“非主”表示它不在所有单位上恒等于一；$\chi(-1)=1$ 通常称为偶特征。要排除它，只需在 $t$ 以内找到一个与 $n$ 互素的整数 $a$，使 $\chi(a)\ne1$。

这不只是二次剩余问题。二次特征仅取 $\pm1$，但一般商群可能有三阶、五阶等特征。只检查平方剩余会漏掉别的乘法障碍。非单位处的零值也不算找到非剩余；准确要求是 $\chi(a)\notin\{0,1\}$。\src{S16}\src{S30}

\section{当总位置数是素数，有一个鸽巢证明}
\begin{proposition}
若 $n=t^2+t+1$ 是素数，则 $H_t=(\Z/n\Z)^\times$。
\end{proposition}
\begin{proof}
给任意非零余数 $a$，考虑 $(t+1)^2$ 个数
\[ai+j\pmod n,\qquad0\le i,j\le t.
\]
它们多于 $n$，所以两项相同。相减得 $a(i-i')=j'-j$。两个差都不为零：若一个为零，由 $n$ 为素数及差的绝对值小于 $n$，另一个也为零，与两项不同矛盾。于是
\[a\equiv\frac{j'-j}{i-i'}\pmod n,
\qquad0<|i-i'|,|j'-j|\le t.
\]
分子分母及符号都在 $H_t$ 的生成元中。
\end{proof}

合数模数时，这个证明的分母可能不是单位。一般“平方根以内的小单位总能生成全部单位群”确实是假的。例如模105时，不超过 $\sqrt{105}$ 的单位只有 $1,2,4,8$；加上负一生成的群有24个元素，而全部单位有48个。105并不是这里的 $t^2+t+1$，所以这反驳的是泛化的证明策略，并不反驳平方族猜想。

\section{无条件文献已经说明大参数一定成功}
Pollack的2017年定理给出：任取 $\varepsilon>0$，存在统一的 $m_0(\varepsilon)$ 和 $\kappa>0$，对每个 $m>m_0$ 及每个非主特征模 $m$，在
\[m^{1/(4\sqrt e)+\varepsilon}
\]
以内有多于 $m^\kappa$ 个素数 $\ell$ 满足 $\chi(\ell)\notin\{0,1\}$。这里只需要其中一个。该定理允许合数模数及非本原特征，符合我们的输入。\cite{Pollack2017}

取 $\varepsilon=1/4-1/(4\sqrt e)>0$，界变成 $n^{1/4}$。对 $t\ge2$，有 $n^{1/4}<t$，所以每个充分大的 $t$ 都满足 $H_t$ 满群。

这已经是无条件结果。它留下的困难很具体：现有引用没有提供足够可用的明确起点 $T_0$。因此，验证 $t<1000$，再说“充分大时成立”，中间仍可能有一段没有覆盖；两句话不能直接拼成全参数定理。

\section{GRH提供一个可计算的交界点}
在Dirichlet $L$ 函数的GRH下，Bach的小非剩余界可写为：每个非主特征模 $n$，都有
\[a\le3(\log n)^2,\qquad\chi(a)\notin\{0,1\}.
\]
项目采用的是Martin--Pollack论文首页对这一界的明确陈述；Bach原论文列在参考文献中，本轮未取得并重读其全文。\cite{MartinPollack2013,Bach1990}

当 $t\ge1000$，有
\[3\bigl(\log(t^2+t+1)\bigr)^2<t.
\]
在 $t=1000$，左边小于573。若把左边记为 $f(t)$，则
\[f'(t)=6\log(t^2+t+1)\frac{2t+1}{t^2+t+1}
<\frac{24\log(t+1)}t<1\qquad(t\ge1000).
\]
所以右边 $t$ 此后增长更快，不等式持续成立。于是GRH给出全部 $t\ge1000$ 的统一覆盖。

剩余 $2\le t<1000$ 共998个参数已经有精确有限证书。其中188个使 $n$ 为素数，可以用上一节证明；其余810个模数使用完整单位群生成核验。二者与GRH范围合起来，得到
\begin{theorem}[平方族的条件性完成]
假设Dirichlet $L$ 函数的GRH。对每个 $t\ge2$，$k=t^2$、$n=t^2+t+1$ 的有限平均条件恰为 $G=1$。
\end{theorem}
这些998个小参数的结论本身是无条件的。对任何使 $t^2+t+1$ 为素数的参数，也有无条件结论，无需限制 $t$ 的大小。

\section{有限证书究竟验证什么}
模 $n$ 的单位群 $\mathcal U$ 先通过CRT分解成素数幂模数的单位群，再写成循环群的直积。每个小整数由一个指数向量表示。要判断这些向量是否生成整个群，可以对每个整除群阶的素数 $\ell$，检查它们在
\[\mathcal U/\mathcal U^\ell
\]
这个有限维 $\F_\ell$ 向量空间中是否张成。

理由很短：若生成子群真小，非平凡有限商群必有一个素数阶商；在对应的 $\mathcal U/\mathcal U^\ell$ 中就不可能满秩。反之，只要每个素数层都满秩，商群就只能平凡。

因此证书核验的是每个指定有限群的完整生成性。它不靠抽样单位，也不把近千个成功参数外推到所有 $t$。源包保存原证书、核验程序及文献访问记录。

\section{现有显式Burgess结果为何还没有直接去掉GRH}
一般短特征和估计当然与此有关，但用于全参数证明时必须同时核对模数条件、特征条件、常数和阈值。项目调查的一份2025年显式结果具有很大的适用起点，并有本原特征或无立方因子等限制。这里 $t^2+t+1$ 不总无立方因子：$t=18$ 时它等于 $343=7^3$。

该来源及其适用版本见\cite{ExplicitBurgess2025}；详细访问等级保存在来源S30中。

所以当前“去掉GRH”的准确任务，是给 $H_t$ 一个覆盖所有参数的无条件证明，或者给一个明确且可衔接的无条件阈值。平方族的物理返回不需要重做；另一方面，现有文献的渐近界也尚未完成这最后的全参数衔接。


\chapter{二、三的奇次幂：同一结构的另外两个周期}\label{ch:twothree}
\section{迹与行列式先告诉我们该试什么}
对一般三块核心 $(u^k,v^r,w)$，$w=-ku-rv$，平均 $u^{k-r},v^{r-1},w$ 给出整数矩阵
\[C=\begin{pmatrix}-r-1&1\\-n&1\end{pmatrix},
\qquad n=k+r+1.
\]
它满足
\[\operatorname{tr}C=-r,\quad\det C=k,\quad C^2+rC+kI=0.
\]
忽略共同尺度后，控制特征值之比的是
\[\theta=\frac{r^2}{k}.
\]
若该比值为单位根，$C$ 才有有限射影阶。因为 $2+2\cos\phi=\theta$ 是有理数，且单位根给代数整数，有理的 $2\cos\phi$ 只能是 $-2,-1,0,1,2$。在 $r>0$ 且矩阵非标量的本模板内，真正的有限阶恰好对应
\[\theta=1,2,3,
\]
射影阶分别为三、四、六。$\theta=4$ 是非平凡Jordan情形，并不产生有限阶。\src{S21}

平方族给 $\theta=1$。若 $k=ms^2,r=ms$，则 $\theta=m$，所以 $m=2,3$ 正好提供另外两类周期。这一计算指导我们选取模板，但不能把 $\theta>4$ 解释成平均问题不可能可解；它只说明这个 $C$ 不再自行循环。

\section{相同的两步选点，新的库存比例}
以下令 $m=2$ 或三，$s\ge2$，并取
\[k=ms^2,\qquad r=ms,\qquad n=ms^2+ms+1.
\]
有
\[C^4=-k^2I\quad(m=2),\qquad C^6=-k^3I\quad(m=3).
\]
因此实际返回 $C/k$ 可逆。另两个一步返回是
\[
\mathcal A=\begin{pmatrix}1&-1/s\\0&-1/s\end{pmatrix},\qquad
\mathcal B=\frac1k\begin{pmatrix}-r-1&r\\-n&k+r\end{pmatrix}.
\]
由 $C^{-1}(k\mathcal B)=\Sigma$，得到 $\Sigma$；再由 $\Sigma\mathcal A=\Delta(1/s)$ 开始建立倍率资源。\src{S17}

对每个 $1\le d\le r$，先选
\[(u^{k-r+d-1},v^{r-d},w)
\]
平均。保留 $r$ 份新值和一份未选的旧 $u$，再平均其余 $k$ 个位置。两步矩阵的整数部分为
\[
M_d=\begin{pmatrix}
r^2+r-k&-rd\\rn&-(r+k)d
\end{pmatrix}=C^2\Delta(d),
\]
实际尺度为 $k^{-2}$。因为 $C$ 有正向逆，直接得到全部 $\Delta(d)$。此处最大 $d=r$ 仍留一份旧 $u$，最小 $d=1$ 也不超出 $v$ 的库存，所以整个参数区间都是实际可用的。

\section{恢复逆元为什么分成几个小情况}
与平方族一样，初始倍率 $s$ 和 $\Sigma$ 只给带因子 $s-1$ 的上剪切，以及 $R_0=\Z[1/s]$ 上的深层。要消掉固定因子，仍希望得到二的可逆倍率。

先补 $m$ 的逆。如果 $m\mid s$，已有的 $\Delta(s)^{-1}\Delta(s/m)$ 直接给出 $\Delta(1/m)$。如果 $m\nmid s(s-1)$，倍率与旧水平互素，有限商方法适用。剩下 $m\mid s-1$ 时，构造
\[
H_m=\frac1{mk}\Delta(m)C^2\Delta(m)
=\begin{pmatrix}
((m-1)s+1)/(ms)&-1/s\\
n/s&-m(s+1)/s
\end{pmatrix}\in\SL_2(R_0).
\]
左上角为 $R_0$ 中元素，正是因为 $s\equiv1\pmod m$。前篇的深层工具给 $H_m$ 的逆，把这个“夹心”式子解开就能得到 $\Delta(1/m)$。

当 $m=3$、$s$ 为奇数时，还需要另补除二。相应行列式一矩阵为
\[
H_2=\frac1{2k^2}\Delta(2)C^4\Delta(2)
=\begin{pmatrix}
(s+1)/(2s)&-1/s\\n/s&-4-2/s
\end{pmatrix}.
\]
它的分母二同样被 $s+1$ 的偶性消去。于是各个分支都得到 $\Delta(2)^{\pm1}$，重算交换子给 $U(1/2)$。

令 $R=\Z[1/(2ms)]$，最终资源与平方族相同：
\[U(R),\quad L(nR),\quad\text{安全小倍率及其逆}.
\]
下剪切来自 $U(-1)(C/k)\Delta(k)=L(-n/k)$。整个推导对所有指定 $m,s$ 都成立。

\section{入口和小单位问题也能接上}
任意输入的入口仍用见证转移。不过这次有更方便的恒等式：
\[m(k-1)-(r-2)(r+2)=4-m.
\]
结合 $n=k+r+1$ 可得 $\gcd(k-1,n)=1$：$m=3$ 时公因子整除一；$m=2$ 时整除二，而 $n$ 为奇数。因此三角色交换子在每个危险素数处都有效。保护位置数量由 $\omega(n)\le r$ 控制，入口成立。

剩下单位群由 $-1$ 和 $1\le d\le r$ 的安全小整数生成。GRH下的小非剩余界在 $s\ge256$ 已落入 $r=ms$ 以内；$2\le s<256$ 的两族共508个参数已有精确有限生成证书。它们接上通用控制器和终端，得到：
\begin{theorem}[二、三平方倍数族]
在Dirichlet $L$ 函数的GRH下，对 $m\in\{2,3\}$、$s\ge2$，$k=ms^2$ 平均在 $n=ms^2+ms+1$ 个位置上恰以 $G=1$ 判定。
\end{theorem}
小参数证书所覆盖的每个具体维数是无条件的。十二平均十九元、十八平均二十五元、五十平均六十一元都属于这两个族；但这些维数未必是各自最早成功点。

\section{为什么这完成了二、三的所有幂的最早成功点}
当 $k=p^a$，临界尺度为 $S(k)=p^{\lceil a/2\rceil}$。偶指数写成平方，上一讲给 $k+S(k)+1$；奇指数且 $p=2,3$ 时写成 $ps^2$，本讲给同一位置数。结合严格下界，对于 $a\ge2$，在GRH下有
\[M(p^a)=H(p^a)=p^a+p^{\lceil a/2\rceil}+1,
\qquad p=2,3.
\]
素数 $k=3$ 由前篇处理；$k=2$ 是 $M=4,H=5$ 的特殊自动点。

这仍未给出 $N(p^a)$。例如九平均十三元成功，没有顺带证明十四元。一个核心解决一个维数；把后续全部维数连接起来，是另一项独立任务。


\chapter{较大素数的奇次幂：已经做出了什么，还缺什么}\label{ch:high}
\section{旧周期消失，不等于问题无解}
现在令 $k=ps^2,r=ps,n=k+r+1$，其中 $p\ge5$。若要研究 $p^{2a+1}$，就取 $s=p^a$。上一讲的基本返回有 $\theta=p>4$，没有有限射影周期。

继续研究的方向是换返回词。当前已找到两类三步返回，其整数矩阵迹为零，因此每个返回重复两次都是标量。它们给出了通用的正向可逆资源，但尚未产生完整剪切群。下面保留其构造，因为它准确展示了进展和剩余的容量问题。\src{S20}

显示的计数与矩阵恒等式实际上允许任意整数 $p\ge5$；本讲关心的应用是 $p$ 为素数的情形。

\section{用三次平均制造一个反射}
假设
\[s\ge(p-1)(p-2).
\]
从核心 $K=(u^k,v^r,w)$ 出发，先取 $j=s-1$，平均
\[(u^{k-j-1},v^j,w).
\]
得到 $k$ 份 $D$，外面剩 $s$ 份旧 $u$ 和 $d=(p-1)s+1$ 份旧 $v$。第二步取 $a$ 份旧 $u$、$b$ 份旧 $v$ 和 $k-a-b$ 份 $D$，得到 $k$ 份 $E$。

第三步保留 $r$ 份 $E$ 和一个指定旧单点，平均其余 $k$ 项。两种可用的第二步计数为
\[
\begin{array}{c|cc}
\text{最终保留}&a&b\\\hline
u&s-(p-1)(p-3)&(p-1)s+3-p\\
v&s-(p-1)(p-2)&(p-1)s+2-p
\end{array}
\]
条件 $s\ge(p-1)(p-2)$ 确保这些数非负；同时它们不超过第一步剩下的库存，而且仍留有所要求的旧单点。这些不等式正是三步词能在原位置上实现的依据。

\section{矩阵公式与它真正证明的范围}
第一行计数对应 $h=p(s-p+3)$，第二行对应 $h=p(s-p+2)$。令
\[\alpha=nh-k(r+1),\qquad\beta=k(d-b)-dh.
\]
三步作用为 $J/k^3$，其中
\[
J=\begin{pmatrix}
-r\alpha-k^2&-r\beta+z\\
-(r+k)\alpha-k^2&-(r+k)\beta+z
\end{pmatrix},
\]
保留 $u$ 时 $z=0$，保留 $v$ 时 $z=k^2$。代入两行计数都得到 $\operatorname{tr}J=0$，所以
\[J^2=\lambda I,\qquad \lambda\equiv k^6\pmod n.
\]
由于 $\gcd(k,n)=1$，$\lambda$ 非零且模 $n$ 可逆。于是两类三步返回各自有正向射影逆。这是对整族参数的恒等式证明。

对真正的奇素数幂，若 $s=p^a$ 且 $a\ge2$，容量条件成立，所以覆盖指数 $2a+1\ge5$。但立方对应 $s=p$，对 $p\ge5$ 不满足这个条件。比如 $p=5$ 要求 $s\ge12$，而立方只有 $s=5$。

\section{两块尚未补齐的内容}
第一，三次幂需要新的计数方案，或允许改变核心形状。现有表中出现负数就意味着没有那么多可选的位置，不能把它视作形式矩阵上的小问题。

第二，即使已经有两个可逆反射，也还要证明它们和其他返回能产生足够的上、下剪切与倍率。两个反射可以生成很小的群；可逆性本身不保证同余群覆盖，更不保证每个合法列都能到终端。

因此当前准确状态是：高奇次幂有两种统一的可逆三步资源；完整核心控制、全输入收尾和精确临界尚未完成。前面的GRH小单位定理不能补这一步，因为它只能判断\emph{已经实际供应的}倍率是否足够。


\chapter{端点已统一，为什么整个区间仍未统一}\label{ch:intervals}
\section{现在已经知道每个元数的2k+1端点}
\begin{theorem}[所有元数的共同端点]
对任意整数 $k\ge2$，在 $n=2k+1$ 个位置上，有限 $k$ 平均可达性恰由 $G=1$ 判定。
\end{theorem}
因为 $\gcd(k,2k+1)=1$，一般的允许素因子条件在这里就是 $G=1$。奇数 $k$ 的证明来自前篇奇余量构造，其实际前提不需要 $k$ 为素数；偶数 $k$ 的补齐使用下面的返回。全部端点定理是无条件的。\src{S18}\src{S19}

这个结论给 $M(k)\le2k+1$，也给 $H(k)\le2k+1$。它不等于 $N(k)\le2k+1$，因为定义 $N$ 时还要求所有后续维数成功。

\section{偶数端点：先取得一个反射，再补除二}
写 $k=2b,n=4b+1$。在两个等大块和一个单点的核心上，已知两步返回的整数部分为
\[T_j=\begin{pmatrix}-k-1&j\\-n&2j\end{pmatrix},
\quad1\le j\le k,\qquad\det T_j=-j.
\]
实际矩阵是 $T_j/k$。保留单点的另一类两步返回为
\[F_h=\begin{pmatrix}1&(h-k)/(2k)\\0&h/k\end{pmatrix},
\qquad h\text{ 为偶数},\quad|h|\le k.
\]
两块等大，所以这次 $\Sigma=\left(\begin{smallmatrix}1&-1\\0&-1\end{smallmatrix}\right)$ 可以直接通过重新登记块角色实现。

令 $J=T_b\Sigma$，有
\[J^2=-bI,\qquad T_b^{-1}T_j=\Delta(j/b).
\]
若 $b$ 为偶数，取 $j=2b,b/2$ 就得到 $\Delta(2)$ 和 $\Delta(1/2)$。于是熟悉的交换子给上剪切，而
\[U(-1/2)T_1\Delta(-1/4)=-\tfrac12L(2n)
\]
给下剪切。二进同余控制器已经就绪。

若 $b\ge3$ 为奇数，先选 $\epsilon\in\{-1,1\}$ 使
\[a_0=(4b+1+\epsilon b)/4
\]
为整数：$b\equiv1\pmod4$ 取负号，$b\equiv3\pmod4$ 取正号。则 $1\le a_0\le k$，它模 $n$ 安全，而且
\[\operatorname{tr}(T_{a_0}F_{2\epsilon})=0.
\]
这个周期给 $F_2$ 的逆。再由
\[\Delta(1/b)F_2^{-1}=U((b-1)/2)
\]
和 $\Sigma$ 对上剪切的变号作用，同时取得 $\Delta(b)^{\pm1}$。初始深层随之建立在 $\Z[1/b]$ 上。

剩下的除二，与平方族采取同一思路。置
\[
J'=\Sigma J\Sigma=\begin{pmatrix}2b&-b\\n&-2b\end{pmatrix},
\quad H=\tfrac12\Delta(2)J'\Delta(2)
=\begin{pmatrix}b&-b\\n&-4b\end{pmatrix}.
\]
$H$ 和 $J'$ 的行列式都为 $b$，所以 $H(J')^{-1}$ 是旧环上的行列式一矩阵。用已有深层取逆，再由
\[J'\Delta(2)H^{-1}=2\Delta(1/2)
\]
取得除二。这样两个奇偶分支接到同一套 $U(\Z[1/2])$、$L(n\Z[1/2])$ 资源。$b=1$ 是前篇的二平均情形。

\section{端点为什么不需要GRH的小生成元定理}
从 $T_b^{-1}T_j$ 可以激活所有 $1\le j\le k$ 的安全倍率及其逆。模 $2k+1$ 的每个非零余数都有绝对值不超过 $k$ 的平衡代表；若余数为单位，这个代表也与模数互素。

因此这里的可用倍率已经覆盖全部单位群。平方族只供应到约 $\sqrt n$，需要小非剩余定理；端点供应到约 $n/2$，平衡代表就足够了。差别来自实际返回能够容纳的计数范围。

贝祖输送把核心送到 $(0^k,(-1)^k,k)$。先平均单点 $k$ 与 $k-1$ 个零，得到 $k$ 份一和一个剩余零；再两次各选 $k/2$ 份一和 $k/2$ 份负一，全体归零。

任意输入入口复用三角色见证转移。此处 $\gcd(k-1,2k+1)$ 至多含三，一个保留单点足够保护这个失效素数，所以偶数性不妨碍入口。

\section{奇合数的部分中间带也已经接通}
\begin{theorem}[奇元数的互素奇余量]
设 $k\ge3$ 为奇数，$1\le r<k$ 为奇数且 $\gcd(k,r)=1$。则 $n=2k+r$ 满足完整 $G=1$ 判据。
\end{theorem}
原素数证明在这部分实际使用的是 $k$ 模危险素数可逆、奇数配比和足够的选点容量。核对短区间筛选后，这些条件对奇合数也成立。入口得到
\[(x^k,(rz-x)^k,(-kz)^r).
\]
两个已知两步反射写成
\[J=\begin{pmatrix}0&d\\-1&0\end{pmatrix},\quad
Q=\begin{pmatrix}0&e\\-1&0\end{pmatrix},\quad
d=-r(k+r)/2,\quad e=d+n.
\]
$\gcd(k,r)=1$ 保证 $d,e$ 互素且均与 $n$ 互素，二者奇偶相反。它们给 $n$ 倍上、下剪切，再由前篇的第三方向与有限环消元取得所需同余层。

最后的小代表也仍可使用：一个单位的平衡代表 $u$ 若大于 $k$，以 $2u-\operatorname{sgn}(u)n$ 替换，绝对值变成小于 $r<k$；再除掉二的幂，得到可由真实数字返回实现的安全奇数。前篇的列准备与终端随后完成。

\begin{corollary}[奇素数幂的全部奇中间维数]
若 $k=p^a$，$p$ 为奇素数，则每个奇数 $n$ 满足 $2k<n<3k$ 时，完整条件都是 $\rad(G)\mid p$。
\end{corollary}
\begin{proof}
令 $r=n-2k$，$d=\gcd(k,r)$，$q=k/d$，$r_0=r/d$。商层有 $m=n/d=2q+r_0$ 个位置，$q,r_0$ 为互素奇数，因此上一条定理适用。

按同时好划分把原位置分成大小 $d$ 的块。因为 $d$ 是 $p$ 幂，商层仍只允许同一分母素数。商层每次选 $q$ 块，相当于原位置选 $qd=k$ 项平均，块均值精确模拟商操作。均值达到目标后，再以覆盖全部块的操作清除可能剩下的块内差异。
\end{proof}
例如九平均二十一、二十三、二十五元都已覆盖。这项结论没有包括偶数位置数。

\section{有公因子时，核心坐标必须先改正}
一般 $n=2k+r$，写 $d=\gcd(k,r)$、$k_0=k/d$、$r_0=r/d$。三块整数零和条件只保证载体值被 $k_0$ 整除，正确参数化为
\[(a^k,(r_0z-a)^k,(-k_0z)^r).
\]
三个块值的最大公因数是 $\gcd(a,z)$。本原时令 $y=a+k_0z$，两个独立差分为 $y$ 和 $2y-(n/d)z$，所以
\[G=\gcd(y,n/d).
\]
例如六平均、$r=3$ 时，核心 $(1^6,0^6,(-2)^3)$ 对应整数参数；若硬写载体为 $-6z$ 并要求 $z$ 为整数，就漏掉了它。\src{S15}

真正的危险模数只由 $n$ 中不整除 $k$ 的素数幂组成。$d$ 的素因子都允许出现在分母里，因此除去 $d$ 不改变危险部分。但在模整个 $n$ 下直接除以 $k$ 仍然不合法。

\section{共同入口不等于共同收尾}
对所有 $k\ge2,n\ge2k+1$，已经有移动见证、再合并的入口，得到两个 $k$ 重块与剩余 $r=n-2k$ 个单点。\src{S26}

它先保护使交换子失效的少数危险素数，再把其他见证通过CRT集中到一个角色，最后在保护位置外做平均。这解决了建立两个大块的问题。它并未自动把 $r$ 个单点压成一个等值块，也未证明之后总能继续降低核心维数。

若要得到 $N(k)\le2k+1$，还必须解决这些剩余参数与更多中间带。更进一步要得到 $N(k)=B(k)$，则还要填补最早成功点至 $2k+1$ 之间的范围。端点、部分奇中间带和大维数尾部是三块已建好的桥面，中间仍有未接通的区间。


\chapter{其他路线、有效启发与已经看清的障碍}\label{ch:obstacles}
\section{允许复制以后，G条件已经是完整答案}
有一个很强、却必须保留“复制”二字的定理。它最早在三平均研究中整理，证明对任意元数都适用。\src{S28}

\begin{theorem}[复制稳定化]
对非零本原整数零和状态 $X$，以下三件事等价：
\begin{enumerate}
\item $G(X)$ 的所有素因子都整除 $k$；
\item 存在 $R=k^h$，将每个坐标复制 $R$ 份后，能够有限步 $k$ 平均归零；
\item 存在条目属于 $\Z[1/k]$ 的严格正双随机矩阵 $P$，使 $PX=0$。
\end{enumerate}
\end{theorem}
\begin{proof}
先由差分的贝祖等式找整数列 $a$，使 $\mathbf1^Ta=0$、$X^Ta=G$。例如将 $G$ 写成各 $x_i-x_1$ 的整数线性组合即可。置
\[b=Ge_1-x_1a,
\qquad\mathbf1^Tb=G,\qquad X^Tb=0.
\]
取充分大的 $R=k^h$ 使 $G\mid R$，写 $R=nq+r$、$0\le r<n$。因为 $G\mid n$，故 $G\mid r$。令
\[m=q\mathbf1+(r/G)b,\qquad
\ell=R\mathbf1-(n-1)m.
\]
当 $q$ 足够大，两列的所有项都是正整数。它们各自的坐标和为 $R$，且与 $X$ 内积都为零。

取矩阵 $M$ 的前 $n-1$ 行为 $m^T$、最后一行为 $\ell^T$。每行和、每列和都是 $R$，且 $MX=0$。所以 $P=M/R$ 是第三项的矩阵；它的秩还至多为二。

矩阵 $M$ 也直接给复制实现：每个原坐标现有 $R$ 份副本，第 $i$ 组取 $M_{ij}$ 份第 $j$ 种值。列和保证恰好用完副本，行和保证每组有 $R=k^h$ 项，$MX=0$ 保证每组和为零。每组用标准 $k$ 进网络平均，即得第二项。

反过来，若某个 $\ell\nmid k$ 整除 $G$，原输入模 $\ell$ 全同余且非零。复制后的任何平均仍保留这个共同余数；而 $\Z[1/k]$ 上的双随机矩阵也因行和为一保留它。因此第二、三项都不可能。
\end{proof}

它说明 $G$ 已完整刻画了一个放宽后的运输问题。原问题还要求同样的消去发生在原来的有限位置上。正双随机矩阵虽写得很漂亮，却未必能在原位置分解为全等权的 $k$ 平均。

24平均27元是这一差别的压力测试：原位置有已证的 $G=1$ 反例，充分复制后却一定可解。因此不能普遍消去副本，也不能把任意非负运输表直接当作原位置平均词。

\section{二点T变换为何不能直接代替我们的原子}
经典主序理论将许多双随机作用分解为两点混合
\[\begin{pmatrix}\lambda&1-\lambda\\1-\lambda&\lambda\end{pmatrix},
\qquad0\le\lambda\le1.
\]
它允许连续的参数 $\lambda$。我们的操作每次必须把恰好 $k$ 个位置全部替换成同一个均值，既没有任意两点参数，也不能随意保留这些位置上的差异。

因此若要借主序理论，真正需要补的是一个模拟引理：用现有位置和允许的平均精确实现所需混合，并且能连续串联。第\ref{ch:transfer}讲的等大块均值模拟已经给出一种有效情形；一般加权运输仍缺这个分解。

\section{单点交换有能力，也有全长度限制}
块与单点交换能转移有限模见证，但未必能在有理数上提供正向逆。项目对素数平均元数 $p$ 的单大块、至少两个单点系统证明了一个刚性结论。\src{S27}

其差分矩阵为
\[A_i=I-e_ie_i^T-\frac1p\mathbf1e_i^T.
\]
每个 $A_i^{-1}$ 都是整数矩阵，$v_p(\det A_i)=-1$。对长度 $L$ 的非空词 $W$，有
\[p^LW\equiv(-1)^L\mathbf1e_{i_1}^T\pmod p.
\]
这控制了Smith标准形的赋值：在 $m\ge2$ 个差分坐标上，指数为 $(-L,0,\ldots,0)$。一个非零标量矩阵的全部Smith指数应相同，所以这里不可能出现非空标量周期。

特别地，这套单点交换字母表内不存在某个生成元的正向射影逆，否则将它接在该生成元后面就得到标量词。这是所有长度的障碍，不是某个搜索深度没找到。

该结论的原证明针对素数载体。这里引用它来解释一种受限方法为何会卡住，没有把它未经证明地推广到所有合数。一般任务中仍可加入多单点操作、改变核心，或只在有限模层恢复逆；后者正是平方入口所用的方法。

\section{每次平均核很大，也不保证当前参数维数下降}
一次 $k$ 平均消掉选中位置的 $k-1$ 维差分空间。不过若当前状态族已经落在一个线性子空间 $L$ 中，真正降低的维数是
\[\dim(L\cap\ker T),
\]
不是整个 $\dim\ker T=k-1$。核很大，却可能与当前像只交于零。

这解释了为何“每次多并几个单点，总会压到二维”的直觉尚不足以构成证明。必须找到与当前参数族有非零交的实际消去方向，并保证同余合法性。一般移动见证入口只解决了前半段资源准备。

同理，原子差分格 $A_{k-1}$ 的秩为 $k-1$，不意味着证明必须进入 $\mathrm{PGL}_{k-1}$。三种块值加一个零和关系仍只有两个参数，无论 $k$ 多大。高秩算术几何可以提示结构，但参数维数应从当前实际核心计算。\src{S24}

\section{逆向设计零和组：一个可保留的具体族}
另一条路线是先规定第二步出现零和组，再倒推第一次的参数条件。它不替代一般判据，但能给一整族显式短词。\src{S22}

例如奇数 $k\ge5$，在 $n=2k+4$ 个位置取
\[(a^k,b^k,c^2,d^2),
\]
其中
\[
c=\frac{(k^2-2k+8)a-(k^2-6k+6)b}{2},
\]
\[
d=\frac{(-k^2+k-8)a+(k^2-7k+6)b}{2}.
\]
它总和为零。先平均 $(a^{k-4},b^3,c)$，得到 $k$ 份
\[z=\frac{ka+(6-k)b}{2}.
\]
随后平均 $(d,c,z,b^{k-3})$，其和恰为零。所有位置库存都足够；两步后有 $k$ 个零，非零支持至多 $k+4\le2k$，可接既有相反块终端程序，整族至多七步。

这给出了一个二维输入族，不是该维数所有 $G$ 合法输入的判据。它说明“先设计终端，再找能进入它的族”很有用，但还需要一个覆盖全部输入的入口。

同一研究中，两个不交碰撞组的保护与合并也提供了通用组合工具：若一次混合撞上旧重值，另一次可保留关键位置，使两个重值的库存同时增加。非单位差分的中点例外则已由第\ref{ch:tail}讲处理。这些局部工具可复用，不等于区间接续已经完成。

\section{哪些旧工作被后续结果替代了}
旧的36个CRT水平曾逐一编译端点群回路，留下32112条精确回路核验。它们说明特定水平能实现，也保留了短词数据。如今全部 $2k+1$ 端点已有一般证明，不能再把这36个水平以外都列为端点未解；旧证书的具体构造与计算记录仍保留。\src{S23}

另一份早期报告依据一篇clique gossip论文的摘要，把“$k\mid n$”作为固定网络必要条件。六平均九元的实际四步网络说明这句话不能直接用于本模型。项目未取得该论文完整定理及其全部假设，故撤回的是这项本地引用推论，并非声称外部论文错误。\src{S05}

还有两种需要保留的研究边界。第一，某个固定模板的周期失败不等于原问题失败；第二，某个有限生成群不够大也不等于所有合法状态不可达。可以有状态相关的路径，或换一个核心。后续研究应瞄准这些具体接口，不再把历史中某一个充分条件当成原问题的唯一出口。


\chapter{把完成的部分与真正剩下的任务放回一张图}\label{ch:map}
\section{先记住四种不同的结论}
一是必要性：对任何位置数，可达必有 $\rad(G)\mid k$。二是充分性：在指定 $(k,n)$ 上，这个条件也足够。三是最早成功点 $M,H$：还要有覆盖此前维数的反例。四是最终阈值 $N$：需要以后每个维数都成功。

本篇的主要进展集中在第二、三类，同时已经有第四类的统一线性上界。精确最终阈值仍是开放任务。下面的状态按实际定理范围整理；原文件中写于较早阶段的“未解”，只在未被后续结果覆盖的范围内保留。\src{S02}\src{S03}\src{S25}

\section{一般定理的覆盖表}
\begin{longtable}{p{.28\textwidth}p{.37\textwidth}p{.25\textwidth}}
\toprule 范围&已经得到的结论&证据与条件\\\midrule\endhead
全部 $k\ge2$&$M,N\ge B$，$H\ge C$&赋值陷阱与边界同余障碍\\
全部 $k\ge3$&$n\ge4k-1+d_k(n)$ 时G等价；$N\le4k+\lceil\log_2k\rceil+1$&无条件，整数能量下降\\
全部 $k\ge2$&$n=2k+1$ 时G等价&无条件，真实返回与同余群\\
偶数 $k$&全部 $k$ 倍数位置数&无条件，因子转译\\
全部 $k$&全部 $n=ak,a\ge3$&无条件，因子转译\\
奇 $k$，奇 $r<k$，$(k,r)=1$&$n=2k+r$ 时G等价&无条件，奇余量控制器\\
奇素数幂 $k$&全部奇数 $2k<n<3k$&无条件，商块转译\\
$k=t^2$&$n=t^2+t+1$ 时G等价&全参数依赖GRH；小参数、素数模数和充分大参数有无条件结论\\
$k=2s^2$ 或 $3s^2$&$n=k+ms+1$，其中 $m=2,3$&全参数依赖GRH；已证有限小参数无条件\\
$k=ps^2,p\ge5$，$s\ge(p-1)(p-2)$&两种三步可逆反射&资源定理；尚不是完整G充分性\\
全部 $k,n\ge2k+1$&可进入两个 $k$ 块加单点的系统&入口定理；尚不是完整收尾\\
\bottomrule
\end{longtable}

平方族对混合素因子不一定打到候选下界。例如 $k=36$ 时，$S(36)=3$，候选是四十元；平方控制器给四十三元。后者再统一，也没有解决前者。

\section{一些具体元数：已证值与候选值分开看}
表中“已证”均指具体参数的无条件结果；条件性全族的量词另列在上一表。
\begin{longtable}{rrrrp{.39\textwidth}}
\toprule $k$&$B(k)$&$C(k)$&已证 $M,H$&来源或尚缺内容\\\midrule\endhead
2&4&5&$4,5$&前篇二平均定理\\
4&7&7&$7,7$&平方核心 $t=2$\\
6&9&10&$9,10$&四步网络、模五控制器\\
8&13&13&$13,13$&$m=2,s=2$ 的具体控制器\\
9&13&13&$13,13$&平方核心 $t=3$\\
10&16&17&$16,17$&六步网络、模十七控制器\\
12&16&17&$16,17$&四步网络、模十七控制器\\
14&22&22&未定&已有 $M,H\le29$ 的端点界\\
15&21&21&未定&候选尚未覆盖；三十一元端点已解\\
16&21&21&$21,21$&平方族小参数证书\\
18&22&22&未定&$m=2,s=3$ 给二十五元，非二十二元\\
20&25&26&$25,\ ?$&四步网络达到 $M$ 下界\\
25&31&31&$31,31$&平方族小参数\\
27&37&37&$37,37$&$m=3,s=3$ 小参数\\
30&36&37&$36,\ ?$&四步网络达到 $M$ 下界\\
32&41&41&$41,41$&$m=2,s=4$ 小参数\\
36&40&40&未定&$40\le M\le43$，平方点为43\\
42&49&50&$49,\ ?$&四步网络达到 $M$ 下界\\
49&57&57&$57,57$&平方族小参数\\
54&64&65&未定&六十四元G自动合法，仍需可达证明\\
56&64&65&$64,\ ?$&四步网络达到 $M$ 下界\\
64&73&73&$73,73$&平方族小参数\\
72&77&77&未定&$77\le M\le81$，上界来自固定网络\\
81&91&91&$91,91$&平方族小参数\\
120&125&126&未定&一百二十五元G自动合法，仍需固定网络或等价证明\\
128&145&145&$145,145$&$m=2,s=8$ 小参数\\
\bottomrule
\end{longtable}

六平均的旧表需要特别更新。十三元由全部 $2k+1$ 端点覆盖，十四与二十二元由素数因子转译直接覆盖。因此在 $9\le n<24$ 中，当前这些定理还未覆盖的六平均位置数是
\[11,16,17,19,20,23.
\]
这是证明覆盖的空白名单，不是不可解维数名单。仍有 $9\le N(6)\le24$。四平均对应 $7\le N(4)\le16$，其中十一、十三、十五元尚未由上述成果接齐。

还有一条常见混淆：旧 $B$ 公式在 $2\le k\le200$ 中有17个自动G点，分别为
\[
\begin{gathered}
2\to4,\ 6\to9,\ 10\to16,\ 12\to16,\ 20\to25,\ 30\to36,\\
42\to49,\ 54\to64,\ 56\to64,\ 75\to81,\ 90\to96,\\
110\to121,\ 120\to125,\ 132\to144,\ 138\to162,\\
156\to169,\ 182\to196.
\end{gathered}
\]
这张表只计算 $\rad(B(k))\mid k$。其中哪些真有固定网络，要另查构造；算术自动合法没有替代这一证明。

\section{三个主任务的精确增量}
\paragraph{任务一：较大素数的奇次幂。}
目标是在 $k=p^{2a+1},p\ge5$ 的候选位置数
\[n=k+p^{a+1}+1
\]
证明完整G充分性。已有的三步反射覆盖 $a\ge2$ 的一部分资源；需要补完整根与倍率，立方还需补容量可行的返回。可以改变核心或使用状态相关路径，未要求一定延续同一个反射模板。

\paragraph{任务二：混合素因子的真正临界点。}
目标是一般 $M(k)=B(k)$、$H(k)=C(k)$。单个危险赋值已经给共同下界，但不同允许素数的局部要求必须由一条原位置路径同时实现。$k=14,15,18,36$ 的候选点和 $k=54,120$ 的自动点，是不同类型的测试，而不是同一份模群代码换参数。

固定网络可处理自动点，低秩控制器可处理非自动点，因子转译可处理合适整除族。真正需要补的是它们在临界库存内的共同实现规则。复制稳定化表明放宽位置后答案很简单，也指出有限库存是不能删去的条件。

\paragraph{任务三：从成功点延续到所有后续维数。}
目标是 $N(k)=B(k)$，较弱的中间目标是 $N(k)\le2k+1$。目前已证端点、部分奇中间带、整除族和线性尾部；一般合数的偶数余量、多危险素数载体折叠与后续终止仍有空缺。需要一个能连续降低某个良序量的选择规则，或一个能保证存在终端路径的完整结构定理。

\section{去GRH与路径效率是另外两项任务}
平方族的GRH依赖已落在明确的 $H_t$ 生成命题上。更强的无条件小生成元界、利用 $\Phi_3(t)=t^2+t+1$ 的特殊算术，或扩大当前可执行倍率，都可能推进它。已经证明的平方入口和周期无需再搜索。

有限可达性也不等于短路径。固定网络族有显式操作数；大维数整数下降有关于数值大小的界；同余群编译保证有限词，却尚未统一控制关于输入位长的原子词长。现有样本中较短的词可以用于算法改进，不能冒充最短步定理。

\section{后续阅读与接手方式}
主讲义把新增结构按依赖次序重写；来源附录逐条指出原证明、实验与核验入口。附带的原文快照保留技术表、程序链接和历史限定，查一个具体公式时可以直接回到对应来源。

本轮新增的是这份续篇、可追溯的来源包及例子核对。没有宣布新的全体合数定理，也没有重新独立审完全部历史证书。接手者应先选择上述某个精确任务，再看其已有资源与缺口；不必重新从“寻找一个统一矩阵群”开始。


\appendix

\chapter{原证明、计算证据与本篇的对应}

本附录以稳定来源编号连接讲义与原研究。压缩包的 sources 目录保持原项目相对路径；来源散列见 SOURCE\_MANIFEST.json。正文按教学顺序重写，原始计数表、推导及历史实验仍保留在对应原文中。

技术原稿中的“未解”或“下一步”可能属于历史阶段；本篇第十四讲按后续定理更新其覆盖范围。核验记录列出原核验的证据等级，本轮没有把所有历史程序重新运行一遍。

原验证器的命令应在原项目根目录执行。证据快照便于检查与追溯，少数历史程序仍可能依赖原工作区；本篇PDF和LaTeX编译不依赖它们。

\section{S01：合数平均的临界尺度：最早成立、最终阈值与无条件网络}\label{source:S01}

M/H/N、赋值下界、B/C猜想、四步网络、八平均十三元；历史状态由S02更新。对应本篇第1—3、6、14讲。

\par\smallskip{\small 原路径：\nolinkurl{outputs/general_arity/composite_critical_scale_and_conjecture.md}}\par

核验登记：{\small\nolinkurl{ten-twelve-seventeen-complete}}。原证据等级：{\small\nolinkurl{theorem-certificate}}。

核验登记：{\small\nolinkurl{seven-average-fifteen-complete}}。原证据等级：{\small\nolinkurl{theorem-certificate}}。

核验登记：{\small\nolinkurl{six-average-ten-complete}}。原证据等级：{\small\nolinkurl{theorem-certificate}}。

核验登记：{\small\nolinkurl{composite-critical-scale}}。原证据等级：{\small\nolinkurl{symbolic-check}}。

核验登记：{\small\nolinkurl{eight-average-thirteen-complete}}。原证据等级：{\small\nolinkurl{theorem-certificate}}。

\section{S02：GRH主线三项任务：本轮完成范围与当前剩余}\label{source:S02}

最近一般元数汇合范围；三项尚未完成的任务。对应本篇第1、10—14讲。

\par\smallskip{\small 原路径：\nolinkurl{outputs/general_arity/general_k_three_tasks_progress_20260915.md}}\par

本文件没有单独匹配的核验登记；其书面推导或依赖结果须在原文中核对，不能据此虚构独立程序。

\section{S03：接受GRH后的完整平方族定理与剩余证明任务}\label{source:S03}

条件性平方族、无条件子族、GRH接口与未闭合量词。对应本篇第8—11、14讲。

\par\smallskip{\small 原路径：\nolinkurl{outputs/general_arity/general_k_grh_stage_and_remaining_tasks_20260915.md}}\par

本文件没有单独匹配的核验登记；其书面推导或依赖结果须在原文中核对，不能据此虚构独立程序。

\section{S04：任意平均元数的线性阈值：非单位差分也是可执行资源}\label{source:S04}

非单位差分、中点碰撞、一般EGZ、统一尾部与4/6造零收尾。对应本篇第5讲。

\par\smallskip{\small 原路径：\nolinkurl{outputs/general_arity/arbitrary_arity_linear_threshold.md}}\par

核验登记：{\small\nolinkurl{arbitrary-arity-linear-threshold}}。原证据等级：{\small\nolinkurl{symbolic-check}}。

\section{S05：合数平均：四平均、六平均与原位置副本转译}\label{source:S05}

秩障碍、副本与半块转译、张量网络、全有理输入引理、摘要误用撤回。对应本篇第3—4、13讲。

\par\smallskip{\small 原路径：\nolinkurl{outputs/general_arity/composite_arity_four_six_transfer.md}}\par

核验登记：{\small\nolinkurl{six-average-nine-fixed-network}}。原证据等级：{\small\nolinkurl{theorem-certificate}}。

核验登记：{\small\nolinkurl{composite-four-six-transfer}}。原证据等级：{\small\nolinkurl{symbolic-check}}。

\section{S06：从素数完整证明到合数临界问题：策略、可迁移定理与真实困难}\label{source:S06}

素数因子块均值转译、全部三倍以上维数及合数特有困难。对应本篇第4、14讲。

\par\smallskip{\small 原路径：\nolinkurl{outputs/general_arity/averaging_proof_strategy_and_composite_frontier_20260915.md}}\par

核验登记：{\small\nolinkurl{square-arity-controller-nine-thirteen}}。原证据等级：{\small\nolinkurl{symbolic-check}}。

核验登记：{\small\nolinkurl{prime-factor-block-mean-transfer}}。原证据等级：{\small\nolinkurl{symbolic-check}}。

\section{S07：四平均七元完整判据：\textbackslash{}((4,2,1)\textbackslash{}) 核与模7控制器}\label{source:S07}

4+2+1真实周期和M(4)=H(4)=7。对应本篇第6—7讲。

\par\smallskip{\small 原路径：\nolinkurl{outputs/general_arity/examples/four_average_seven_complete.md}}\par

核验登记：{\small\nolinkurl{four-average-seven-complete}}。原证据等级：{\small\nolinkurl{theorem-certificate}}。

\section{S08：四平均九元完整判据：一个四步周期与模9控制器}\label{source:S08}

偶元数端点种子与模9完整陪集；后由S18推广。对应本篇第6、12讲。

\par\smallskip{\small 原路径：\nolinkurl{outputs/general_arity/examples/four_average_nine_complete.md}}\par

核验登记：{\small\nolinkurl{four-average-nine-complete}}。原证据等级：{\small\nolinkurl{theorem-certificate}}。

\section{S09：六平均九元：所有输入固定四步完成}\label{source:S09}

固定四步网络及固定词三步下界。对应本篇第3讲。

\par\smallskip{\small 原路径：\nolinkurl{outputs/general_arity/examples/six_average_nine_fixed_network.md}}\par

核验登记：{\small\nolinkurl{six-average-nine-fixed-network}}。原证据等级：{\small\nolinkurl{theorem-certificate}}。

\section{S10：六平均十元完整G判据：第一个非自动的算术成功维数}\label{source:S10}

G=1或2、模5控制器、两步返回与H(6)=10。对应本篇第1、6讲。

\par\smallskip{\small 原路径：\nolinkurl{outputs/general_arity/examples/six_average_ten_complete.md}}\par

核验登记：{\small\nolinkurl{six-average-ten-complete}}。原证据等级：{\small\nolinkurl{theorem-certificate}}。

\section{S11：十平均与十二平均17元完整G判据，以及原B公式中的自动维数}\label{source:S11}

两套实际返回、模17尺度、自动B点表、H(10)=H(12)=17。对应本篇第3、6、14讲。

\par\smallskip{\small 原路径：\nolinkurl{outputs/general_arity/ten_twelve_average_seventeen_complete.md}}\par

核验登记：{\small\nolinkurl{ten-twelve-seventeen-complete}}。原证据等级：{\small\nolinkurl{theorem-certificate}}。

\section{S12：十平均十六元与一族确定性三块网络}\label{source:S12}

固定2b+2步网络族、十平均十六元。对应本篇第3讲。

\par\smallskip{\small 原路径：\nolinkurl{outputs/general_arity/ten_average_sixteen_and_dyadic_network_family.md}}\par

核验登记：{\small\nolinkurl{ten-sixteen-dyadic-network-family}}。原证据等级：{\small\nolinkurl{theorem-certificate}}。

\section{S13：平方元数的统一正向控制器与九平均十三元}\label{source:S13}

平方核心全参数周期、两步数字返回、除二与完整根。对应本篇第7讲。

\par\smallskip{\small 原路径：\nolinkurl{outputs/general_arity/square_arity_uniform_controller_and_nine_thirteen.md}}\par

核验登记：{\small\nolinkurl{square-arity-controller-nine-thirteen}}。原证据等级：{\small\nolinkurl{symbolic-check}}。

\section{S14：一般k平均规划评议：平方族入口可直接补齐，尺度任务可以数论化}\label{source:S14}

全t平方入口、H\_t、一步终端尺度限制。对应本篇第8、13讲。

\par\smallskip{\small 原路径：\nolinkurl{outputs/general_arity/general_k_plan_review_and_square_entry_20260915.md}}\par

核验登记：{\small\nolinkurl{square-core-entry-review}}。原证据等级：{\small\nolinkurl{symbolic-check}}。

\section{S15：对一般k规划建议的复核与采纳}\label{source:S15}

Gamma1扩环接口、正确一般整数核心与G公式。对应本篇第8、12讲。

\par\smallskip{\small 原路径：\nolinkurl{outputs/general_arity/general_k_review_response_20260915.md}}\par

核验登记：{\small\nolinkurl{square-core-entry-review}}。原证据等级：{\small\nolinkurl{symbolic-check}}。

\section{S16：H\_t文献调查：已有小生成元理论、无条件充分大结论与全范围边界}\label{source:S16}

小生成元、特征、Pollack、Bach/GRH、998参数证书、105及立方因子反例。对应本篇第9讲。

\par\smallskip{\small 原路径：\nolinkurl{outputs/general_arity/ht_literature_and_eventual_square_completion_20260915.md}}\par

核验登记：{\small\nolinkurl{ht-literature-interfaces}}。原证据等级：{\small\nolinkurl{theorem-certificate}}。

\section{S17：统一四／六阶控制器完成2、3奇次幂临界点}\label{source:S17}

2s²/3s²完整控制器、508组小参数与条件性临界族。对应本篇第10讲。

\par\smallskip{\small 原路径：\nolinkurl{outputs/general_arity/two_three_odd_power_critical_grh_completion.md}}\par

核验登记：{\small\nolinkurl{two-three-odd-power-critical}}。原证据等级：{\small\nolinkurl{theorem-certificate}}。

\section{S18：全部偶数元数的2k+1端点：旧种子与统一取逆}\label{source:S18}

全部偶数元数2k+1端点；接奇数端点得到所有k。对应本篇第12讲。

\par\smallskip{\small 原路径：\nolinkurl{outputs/general_arity/even_arity_all_endpoints_completion.md}}\par

核验登记：{\small\nolinkurl{even-arity-all-endpoints}}。原证据等级：{\small\nolinkurl{symbolic-check}}。

\section{S19：奇合数元数的奇余量区间：按真实互素条件迁移旧证明}\label{source:S19}

奇k奇r互素带、奇素数幂全部奇中间维数。对应本篇第12讲。

\par\smallskip{\small 原路径：\nolinkurl{outputs/general_arity/odd_composite_coprime_middle_band_completion.md}}\par

核验登记：{\small\nolinkurl{odd-composite-coprime-middle-band}}。原证据等级：{\small\nolinkurl{symbolic-check}}。

\section{S20：较大素数高奇次幂：两种统一三步正向反射}\label{source:S20}

p>=5高奇次幂的两种三步反射，未得完整根；立方未覆盖。对应本篇第11讲。

\par\smallskip{\small 原路径：\nolinkurl{outputs/general_arity/large_prime_high_odd_power_three_step_reflections.md}}\par

核验登记：{\small\nolinkurl{large-prime-high-odd-power-reflections}}。原证据等级：{\small\nolinkurl{symbolic-check}}。

\section{S21：一般k平均的证明计划：先完成核心，再接成全部维数}\label{source:S21}

theta周期分类、临界／混合／接续三项任务及验收边界。对应本篇第10—14讲。

\par\smallskip{\small 原路径：\nolinkurl{outputs/general_arity/general_k_averaging_proof_plan_20260915.md}}\par

核验登记：{\small\nolinkurl{square-core-entry-review}}。原证据等级：{\small\nolinkurl{symbolic-check}}。

\section{S22：逆零和定理的原位置转译：碰撞修复与 2p+4 元七步族}\label{source:S22}

通用双碰撞修复及奇数元数七步二参数族；素数加权结果仅引用前史。对应本篇第13讲。

\par\smallskip{\small 原路径：\nolinkurl{outputs/general_arity/inverse_zero_sum_averaging_progress.md}}\par

核验登记：{\small\nolinkurl{inverse-zero-sum-transfer}}。原证据等级：{\small\nolinkurl{symbolic-check}}。

\section{S23：从十五元推广到合数端点：统一CRT证书与两素数完整入口}\label{source:S23}

36个CRT陪集证书的历史阶段及后续替代范围，不重教素数篇陪集算法。对应本篇第12、13讲。

\par\smallskip{\small 原路径：\nolinkurl{outputs/general_arity/composite_endpoint_crt_transfer.md}}\par

核验登记：{\small\nolinkurl{composite-endpoint-transfer}}。原证据等级：{\small\nolinkurl{theorem-certificate}}。

\section{S24：q-平均的根格、\textbackslash{}(S\textbackslash{})-算术与控制器猜想}\label{source:S24}

原子秩与压缩参数秩的区别；几何／正运输仅作为有边界的候选框架。对应本篇第13讲。

\par\smallskip{\small 原路径：\nolinkurl{outputs/general_arity/triple_average_q_average_arithmetic_geometry.md}}\par

本文件没有单独匹配的核验登记；其书面推导或依赖结果须在原文中核对，不能据此虚构独立程序。

\section{S25：从二平均到任意平均元数：当前全貌与临界猜想}\label{source:S25}

历史全局表与最新定理对照；不沿用过时未解项。对应本篇第14讲。

\par\smallskip{\small 原路径：\nolinkurl{outputs/general_arity/averaging_global_status_20260912.md}}\par

本文件没有单独匹配的核验登记；其书面推导或依赖结果须在原文中核对，不能据此虚构独立程序。

\section{S26：移动见证再合并：全部元数、全部超临界维数的同一入口}\label{source:S26}

所有k、n>=2k+1的移动见证双块入口，未提供全范围终止。对应本篇第12—13讲。

\par\smallskip{\small 原路径：\nolinkurl{research/prime_arity_padic/move_witnesses_then_merge.md}}\par

本文件没有单独匹配的核验登记；其书面推导或依赖结果须在原文中核对，不能据此虚构独立程序。

\section{S27：用载体的 p 进格位移审查统一机制}\label{source:S27}

素数单载体全字母表刚性、秩下降交条件、不能凭群或运输代替原位置词。对应本篇第13讲。

\par\smallskip{\small 原路径：\nolinkurl{research/prime_arity_padic/carrier_balance_and_uniform_mechanism.md}}\par

本文件没有单独匹配的核验登记；其书面推导或依赖结果须在原文中核对，不能据此虚构独立程序。

\section{S28：任意维数的新统一结果：复制稳定化、双随机矩阵与消去猜想}\label{source:S28}

第2及8节：任意k复制稳定化、正双随机证书及原位置消去缺口。对应本篇第13讲。

\par\smallskip{\small 原路径：\nolinkurl{outputs/triple_average/history/triple_average_stabilization_theorem.md}}\par

核验登记：{\small\nolinkurl{stabilization-theorem}}。原证据等级：{\small\nolinkurl{finite-sanity}}。

\section{S29：任意维数的三幂分块压缩：把一般输入降到对数秩标准核}\label{source:S29}

同时好划分的独立组合依据。对应本篇第4讲。

\par\smallskip{\small 原路径：\nolinkurl{outputs/triple_average/history/triple_average_power_partition_reduction.md}}\par

核验登记：{\small\nolinkurl{power-partition-compression}}。原证据等级：{\small\nolinkurl{finite-sanity}}。

\section{S30：H\_t文献来源与访问等级}\label{source:S30}

外部文献版本、已读范围和未取得全文的边界。对应本篇第9、来源附录讲。

\par\smallskip{\small 原路径：\nolinkurl{work/ht_literature_20260915/source_index.md}}\par

本文件没有单独匹配的核验登记；其书面推导或依赖结果须在原文中核对，不能据此虚构独立程序。


\backmatter
\begin{thebibliography}{9}
\bibitem{Pollack2017}
Paul Pollack.
\emph{Bounds for the first several prime character nonresidues}.
Proceedings of the American Mathematical Society, 145 (2017), 2815--2826.
\href{https://doi.org/10.1090/proc/13432}{doi:10.1090/proc/13432}.
使用Theorem 1.1的小非剩余界；本轮直接核对本地全文提取文本的该定理。

\bibitem{MartinPollack2013}
Greg Martin and Paul Pollack.
\emph{The average least character non-residue and further variations on a theme of Erdős}.
Journal of the London Mathematical Society, 87 (2013), 22--42.
\href{https://doi.org/10.1112/jlms/jds036}{doi:10.1112/jlms/jds036}.
本篇使用印刷第22页对Bach条件界的陈述，不以平均值结果代替最坏情形结果。

\bibitem{Bach1990}
Eric Bach.
\emph{Explicit bounds for primality testing and related problems}.
Mathematics of Computation, 55 (1990), 355--380.
\href{https://doi.org/10.1090/S0025-5718-1990-1023756-8}{doi:10.1090/S0025-5718-1990-1023756-8}.
此项为经上一条文献核对的间接引用；本轮没有取得并重读原证明。

\bibitem{ExplicitBurgess2025}
Elchin Hasanalizade, Hua Lin, Greg Martin, Andradis Luna Martínez, Enrique Treviño.
\emph{Explicit Burgess inequalities for cubefree moduli}.
\href{https://arxiv.org/abs/2511.17778}{arXiv:2511.17778}, 2025，项目于2026年9月15日保存的作者版本。
仅用于说明现有显式界的适用限制，不作为本篇全参数定理的依据。
\end{thebibliography}

素数篇已解释并核对的EGZ、Morris与Serre工具在本篇按原适用前提复用。其原始题录与具体定理版本仍以素数篇及附带研究原文中的来源记录为准。本篇新增调用的特征非剩余文献列于上方。
\end{document}
